% !TeX program = xelatex
% Self-contained source. Compile twice with XeLaTeX.
% Required fonts: Noto Serif CJK SC and Noto Sans CJK SC.
% Standard TeX Live packages only; no external images, bibliography, or data required.
% Frozen content.json SHA256: 767988f7bddca24a079bb2cb75d655c816a14ff5869d3ebd8eb3843d92272996
\documentclass[11pt,a4paper]{article}
\usepackage{fontspec}
\usepackage[margin=22mm,top=21mm,bottom=22mm,headheight=14pt,headsep=8mm,footskip=10mm]{geometry}
\usepackage{graphicx,array,longtable,booktabs,xcolor,enumitem,needspace,titlesec,fancyhdr}
\usepackage[unicode,bookmarksopen=true,bookmarksnumbered=false]{hyperref}
\usepackage{xurl}
\setmainfont{Noto Serif CJK SC}[BoldFont={Noto Serif CJK SC Bold}]
\setsansfont{Noto Sans CJK SC}[BoldFont={Noto Sans CJK SC Bold}]
\setmonofont{DejaVu Sans Mono}[Scale=0.85]
\XeTeXlinebreaklocale "zh"
\XeTeXlinebreakskip=0pt plus 0.6pt
\XeTeXlinebreakpenalty=0
\definecolor{LinkBlue}{HTML}{1F4E68}
\definecolor{Muted}{HTML}{515B62}
\hypersetup{colorlinks=true,linkcolor=black,urlcolor=LinkBlue,pdftitle={URP研究与实施方案},pdfsubject={农业路径规划算法库 智能推荐与自动算法设计},pdfauthor={}}
\urlstyle{same}
\setcounter{secnumdepth}{0}
\setcounter{tocdepth}{2}
\renewcommand{\contentsname}{目录}
\renewcommand{\tablename}{表}
\setlength{\parindent}{0pt}
\setlength{\parskip}{5.3pt plus 0.5pt minus 0.5pt}
\linespread{1.28}
\emergencystretch=2em
\hyphenpenalty=10000
\tolerance=1800
\hbadness=5000
\widowpenalty=10000
\clubpenalty=10000
\displaywidowpenalty=10000
\raggedbottom
\setlist[itemize]{leftmargin=1.3em,label=\textbullet,itemsep=2pt,topsep=4pt,parsep=0pt}
\titleformat{\section}{\sffamily\bfseries\fontsize{20}{28}\selectfont\color{black}}{}{0pt}{}
\titleformat{\subsection}{\sffamily\bfseries\fontsize{14}{21}\selectfont\color{black}}{}{0pt}{}
\titleformat{\subsubsection}{\sffamily\bfseries\fontsize{11.2}{17}\selectfont\color{black}}{}{0pt}{}
\titlespacing*{\section}{0pt}{0pt}{16pt}
\titlespacing*{\subsection}{0pt}{17pt}{7pt}
\titlespacing*{\subsubsection}{0pt}{11pt}{4pt}
\pagestyle{fancy}
\fancyhf{}
\fancyhead[L]{\sffamily\fontsize{8}{10}\selectfont\color{Muted}URP研究与实施方案}
\fancyhead[R]{\sffamily\fontsize{8}{10}\selectfont\color{Muted}\nouppercase{\leftmark}}
\fancyfoot[C]{\sffamily\fontsize{9}{11}\selectfont\thepage}
\renewcommand{\headrulewidth}{0.3pt}
\renewcommand{\sectionmark}[1]{\markboth{#1}{}}
\setlength{\LTpre}{6pt}
\setlength{\LTpost}{8pt}
\newcommand{\bodytext}[1]{#1\par}
\newcommand{\sourcetext}[1]{{\fontsize{9.2}{14}\selectfont\color{Muted}#1\par}}
\newcommand{\referencetext}[1]{{\fontsize{9.4}{14.2}\selectfont#1\par}}
\begin{document}
\pagenumbering{roman}
\thispagestyle{empty}

% content block 0:1 title
{\sffamily\bfseries\fontsize{27}{36}\selectfont URP研究与实施方案\par}\vspace{9pt}
% content block 0:1 subtitle
{\sffamily\fontsize{14}{22}\selectfont 农业路径规划算法库 智能推荐与自动算法设计\par}\vspace{8pt}
% content block 0:1 meta
{\fontsize{10}{15}\selectfont\color{Muted}面向农业机器人的路径规划算法库构建与智能推荐方法研究\newline 统一方案  2026年10月11日\par}\vspace{14pt}\hrule\vspace{12pt}
% content block 0:1 h2
\Needspace{5\baselineskip}{\sffamily\bfseries\fontsize{12}{18}\selectfont 统一路线\par}\vspace{2pt}
% content block 0:1 p
\bodytext{以同一套可执行农业路径规划系统承接课题与研究：建设按能力注册的算法库，让LLM理解场景、提取约束、提出候选算法或受限执行计划，真实调用求解器并独立验算；离线自动设计产生的程序经准入后进入同一算法库。老师的思想与代码双表示演化、多问题自动设计主线保留，并与在线理解和推荐分别评价。}
% content block 0:2 p
\bodytext{交付须覆盖图搜索、采样规划、局部避障和任务优化四类算法能力，以及果园运输、农田巡检避障、多作业点规划和不规则地块覆盖四类案例。VRP适合重点实证，点到点导航、整行服务和面积覆盖保留各自合同与入口。旧11算法名单可用现代强方法替代，替代须有能力映射和可复查证据。}
% content block 0:2 h2
\Needspace{5\baselineskip}{\sffamily\bfseries\fontsize{12}{18}\selectfont 实施与研究同步推进\par}\vspace{2pt}
% content block 0:3 p
\bodytext{近期先把已有CLI、checker、选择与生成资产接到现有React/\allowbreak{}HTTP面板，以OpenLayers、pyproj、Shapely和版本化任务内核实现编辑、运行、核验与回放。首个点服务或果园合同合格后即可接在线LLM；行服务、补给和覆盖能力并行补齐，不等待全部性能矩阵完成。}
% content block 0:4 p
\bodytext{实验先识别受约束理解和真实推荐的增量，再以强传统库与生成扩展库、固定策略与实际选择的四格比较拆开扩库收益。学习型路线改进、动态预算、代码修复、开放模型后训练、生成多样性、agent控制、世界模型及交互研究均保留最小实验。建议集中单次实验的未知量，不把投入排序写成已取消的方向。}
% content block 0:3 h2
\Needspace{5\baselineskip}{\sffamily\bfseries\fontsize{12}{18}\selectfont 本方案的使用方式\par}\vspace{2pt}
% content block 0:5 p
\bodytext{前半部规定需求、现状、任务与系统实现；中部给出LLM主链和探索实验；后半部统一数据、成本、阶段门与验收。已有结果按保存原件陈述，拟议实验和性能目标单独标明。本文整合三份定稿，未启动新实验、安装或额外文献调查。}
% content block 0:1 source
\sourcetext{依据《URP研究计划》《URP技术方案》《研究路线评议》2026年10月11日定稿及已核NDS候选入口。既有冲突以最新研究计划的已审计事实为准；地图实施以技术方案为准。}
\clearpage
\tableofcontents
\clearpage
\pagenumbering{arabic}
% content block 1:1 h1
\section{阅读导航}
% content block 1:6 p
\bodytext{本方案按需求、证据、实现、实验和验收组织。正文采用统一标题层级，可在Word导航窗格中直接跳转。}
% content block 1:1 table
\begingroup
\fontsize{9.6}{14.6}\selectfont
\setlength{\tabcolsep}{4pt}
\renewcommand{\arraystretch}{1.18}
\begin{longtable}{@{}>{\raggedright\arraybackslash}p{\dimexpr 0.331412\linewidth-5.3026pt\relax}>{\raggedright\arraybackslash}p{\dimexpr 0.668588\linewidth-10.6974pt\relax}@{}}
\toprule
{\sffamily\bfseries 部分} & {\sffamily\bfseries 主要内容} \\
\midrule
\endfirsthead
\toprule
{\sffamily\bfseries 部分} & {\sffamily\bfseries 主要内容} \\
\midrule
\endhead
\midrule
\multicolumn{2}{r}{\fontsize{8}{10}\selectfont\color{Muted}续下页}\\
\endfoot
\bottomrule
\endlastfoot
一 研究要求与现代实现 & 课题边界、能力替代、正式成果 \\
\addlinespace[4pt]
二 已有证据与农业任务 & 保存结果、农业任务、服务与转向资格 \\
\addlinespace[4pt]
三 地图 系统与执行实施 & OpenLayers、坐标、输入、交互、Ubuntu与执行API \\
\addlinespace[4pt]
四 LLM理解 推荐与自动设计 & 在线LLM与离线生成、强基线和扩库四格 \\
\addlinespace[4pt]
五 探索方向与最小实验 & 学习路线、动态预算、修复进化、后训练、多样性、harness、世界模型、交互 \\
\addlinespace[4pt]
六 数据 成本与风险 & 独立母组、两层留出、全成本、失败、风险校准 \\
\addlinespace[4pt]
七 分阶段实施与交付 & 滚动阶段门、论文主张与交付证据 \\
\addlinespace[4pt]
附录 来源覆盖与参考依据 & 阅读覆盖、16主题索引、项目原件及参考链接 \\
\end{longtable}
\endgroup
% content block 1:4 h2
\Needspace{6\baselineskip}\subsection{术语与状态约定}
% content block 1:7 p
\bodytext{任务合同指输入、服务定义、目标、约束和可接受输出的明确约定。算法能力合同声明某版本支持的任务与约束组合。runner为受控执行器，checker为独立验算器，selector为实际选择器。SBS是训练侧选定的固定最佳方法，VBS是事后逐例最优的虚拟上界；后者不能部署为免费推荐。}
% content block 1:8 p
\bodytext{“已核”指来源、代码或保存结果已核对；“已实现”只说明相应代码存在；“当前通过”需本环境实际运行；“建议”“拟定目标”均是待实施。系统仿真、半合成作业与田间验证分开陈述。}
% content block 1:9 p
\bodytext{来源R编号沿用主要研究依据，T编号对应工程文档，ER编号对应探索方法，统一收录于文末；它们是检索标签，不代表研究优先级。}
% content block 2:2 h1
\clearpage
\hypertarget{part_1}{}
\section{一 研究要求与现代实现}
% content block 2:5 h2
\Needspace{6\baselineskip}\subsection{1  研究目标与验收范围}
% content block 2:10 p
\bodytext{正式课题是《面向农业机器人的路径规划算法库构建与智能推荐方法研究》。本方案按最新研究意图解释申报材料：保持研究目标、能力覆盖和交付期望，允许用更合适的算法、接口与实验组织替代旧实现。升级的判断依据是能力是否完整、证据是否更强，而不是是否复刻所有旧名字。}
% content block 2:1 h3
\Needspace{4\baselineskip}\subsubsection{三个来源具有不同约束力}
% content block 2:11 p
\bodytext{立项申报书明确算法库、场景理解与约束提取、候选推荐、对接具体算法得到规划结果，以及五项成果。四类算法与四类农业案例描述了预期能力。文中列名的11算法、Python和API工作流是当时的实现路径，后续可以改进，但需解释替代覆盖了哪些要求。[R1]}
% content block 2:12 p
\bodytext{老师《LLM自动算法设计入门指南》第六章名为“实践项目建议”。其中P3强调思想与代码双表示演化，P5强调LLM、进化算法和评估模块的多问题集成，P6期待原创方法与可投稿论文。这些支持保留真实自动设计研究，但不能将全部学习建议改写为学校逐项硬验收。[R2]}
% content block 2:13 p
\bodytext{后续实施方案将固定车队CVRP、VRPTW、IL与PPO等作为实施取舍。未找到老师逐项批准其替代全部农业能力的原始记录。因此它可帮助接续已有成果，不能用来取消导航、避障或覆盖任务。[R3]}
% content block 2:2 h3
\Needspace{4\baselineskip}\subsubsection{能力替代与交付追踪}
% content block 2:1 bullets
\begin{itemize}
\item 记录原研究需求、拟采用的新实现、替代理由、能力边界和可复查验收证据。
\item 传统方法与学习方法共用任务定义和独立验算；算法名称相同不代表合同相同。
\item 把“已实现”“历史结果已核对”“当前环境已重跑”和“拟研究”分开登记。
\item 最终提交前与导师核对实质性范围变化及学校材料要求；不将没有找到的验收通知补成事实。
\end{itemize}
% content block 2:14 p
\bodytext{申报执行期为一年，研究安排写为2026年3月至2027年2月；尚未找到精确到日的学校截止通知。本方案据依赖和证据安排阶段，不为尚未估算的实验承诺固定周数。[R1]}
% content block 2:6 h2
\Needspace{6\baselineskip}\subsection{2  能力覆盖与现代实现}
% content block 2:15 p
\bodytext{算法库采用按任务能力注册的方式。优先保留已验证的强方法与简明参照，每个成员标明支持的任务、约束、预算与许可证；不同参数配置可以是不同候选，但不包装成新的算法家族。下表是建议实现，不是声称这些能力已经完成。}
% content block 2:2 table
\begingroup
\fontsize{9.6}{14.6}\selectfont
\setlength{\tabcolsep}{4pt}
\renewcommand{\arraystretch}{1.18}
\begin{longtable}{@{}>{\raggedright\arraybackslash}p{\dimexpr 0.155620\linewidth-3.7349pt\relax}>{\raggedright\arraybackslash}p{\dimexpr 0.417867\linewidth-10.0288pt\relax}>{\raggedright\arraybackslash}p{\dimexpr 0.426513\linewidth-10.2363pt\relax}@{}}
\toprule
{\sffamily\bfseries 研究要求} & {\sffamily\bfseries 建议现代实现} & {\sffamily\bfseries 验收证据} \\
\midrule
\endfirsthead
\toprule
{\sffamily\bfseries 研究要求} & {\sffamily\bfseries 建议现代实现} & {\sffamily\bfseries 验收证据} \\
\midrule
\endhead
\midrule
\multicolumn{3}{r}{\fontsize{8}{10}\selectfont\color{Muted}续下页}\\
\endfoot
\bottomrule
\endlastfoot
图搜索与几何可达 & 保留A*/\allowbreak{}Dijkstra精确参照，按车辆状态扩展搜索；不为凑数重复所有旧变体。 & 同图合同、小例参照、可达与不可达、足迹/\allowbreak{}目标核对。 \\
\addlinespace[4pt]
采样与连续运动 & 选定一个匹配运动约束的强采样规划器或组合，保留其支持域。 & 同连续空间与预算，碰撞/\allowbreak{}运动检查、成功率和质量；声明离散精度。 \\
\addlinespace[4pt]
局部避障 & 采用匹配车辆模型的局部规划或控制器；明确动态障碍与全局引导。 & 固定仿真步长和信息权限；到达、碰撞、卡住、控制时延及重放。 \\
\addlinespace[4pt]
任务优化 & PyVRP/\allowbreak{}HGS类强方法、精确小例与合格生成程序；不受GA/\allowbreak{}PSO/\allowbreak{}ACO旧名单限制。 & 任务服务、容量/\allowbreak{}车队/\allowbreak{}时间规则正确；强基线与全分母失败。 \\
\addlinespace[4pt]
行与面积作业 & 保留方向、入出状态与资源；以Fields2Cover等适配不规则覆盖。 & 整行及工具扫掠验算；漏作、重叠、转弯、补给和完整成本。 \\
\addlinespace[4pt]
LLM推荐及调用 & 受约束理解与计划、能力门控、真实性能辅助、注册ID实际求解。 & 原意对齐、LLM提议与执行ID、工具回执、独立验算及实测对照。 \\
\end{longtable}
\endgroup
% content block 2:16 p
\bodytext{现代替代必须保持四类算法能力及四案例的真实覆盖。旧11算法不逐项锁死；每个替代记录输入、输出、支持域、失败和真实调用。能力门控不应成为回避尚缺类别的理由，也不要求不同职责算法进入同一排行榜。}
% content block 2:2 source
\sourcetext{需求依据 R1；方法参考 R8、R12、R13。地图、接口及工程验收在本方案技术部分统一展开。}
% content block 3:3 h1
\clearpage
\hypertarget{part_2}{}
\section{二 已有证据与农业任务}
% content block 3:7 h2
\Needspace{6\baselineskip}\subsection{1  已有资产与可复用结论}
% content block 3:17 p
\bodytext{以下承接2026年10月11日原件核对的保存结果；本次整合未重新求解。已有真实生成、选择和农业组件实验，应复用其代码、协议与失败记录，再补当前环境的运行资格。[R4]}
% content block 3:3 table
\begingroup
\fontsize{9.6}{14.6}\selectfont
\setlength{\tabcolsep}{4pt}
\renewcommand{\arraystretch}{1.18}
\begin{longtable}{@{}>{\raggedright\arraybackslash}p{\dimexpr 0.154179\linewidth-3.7003pt\relax}>{\raggedright\arraybackslash}p{\dimexpr 0.420749\linewidth-10.0980pt\relax}>{\raggedright\arraybackslash}p{\dimexpr 0.425072\linewidth-10.2017pt\relax}@{}}
\toprule
{\sffamily\bfseries 资产} & {\sffamily\bfseries 已核实的结果} & {\sffamily\bfseries 下一阶段用途与边界} \\
\midrule
\endfirsthead
\toprule
{\sffamily\bfseries 资产} & {\sffamily\bfseries 已核实的结果} & {\sffamily\bfseries 下一阶段用途与边界} \\
\midrule
\endhead
\midrule
\multicolumn{3}{r}{\fontsize{8}{10}\selectfont\color{Muted}续下页}\\
\endfoot
\bottomrule
\endlastfoot
BPP选择 & 768次完整执行；训练48母多重集、确认36母多重集，哈希无交集。确认相对训练SBS平均少21.454箱。 & 复用冻结模型、特征、全分母与反例；三种顺序相关，不能当独立母实例。 \\
\addlinespace[4pt]
VRPTW选择 & 40训练母几何、24测试母几何；240训练＋144测试求解记录。8特征48/\allowbreak{}48成功，距离改善1.444\%。 & 复用性能矩阵与部署账；4特征45/\allowbreak{}48成功。只是固定三配置选择，尚非生成新算法闭环。 \\
\addlinespace[4pt]
P2批次LLM生成 & 60槽，43完整评分、7静态拒绝、10生成失败；有提示、父代、回答和真实CLI回执。 & 复用谱系、准入、评分与失败账。EOH臂共享父代快照，不能当独立发现重启。 \\
\addlinespace[4pt]
NL124 & 124母田块、1488槽；固定P为744/\allowbreak{}744可行并达自身矩阵精确参照，默认A为625/\allowbreak{}744。 & 复用精确回归与目标错配诊断。每田12个合成点，不是面积覆盖。 \\
\addlinespace[4pt]
运动与导航 & turn为12根图×3规则；Dubins为24几何×3半径；果园12合成图、96查询。 & 复用模式、姿态、碰撞与连接验算。重复计时不能增加独立地图数。 \\
\end{longtable}
\endgroup
% content block 3:18 p
\bodytext{P102还有24个基础实例在3容量、2空间条件下的144条结果；Fields2Cover源码与350份WKT存在，当前PyVRP适配器max\_\allowbreak{}reloads=0。未找到URP完整行服务、回补、覆盖与运动联合闭环的原始结果。不要把上游仓库能力算作本项目已完成。[R4]}
% content block 3:3 h3
\Needspace{4\baselineskip}\subsubsection{结果级资产的边界}
% content block 3:19 p
\bodytext{NL124为124田块，每块一组仓库与12个单位需求客户，容量5、固定3条非空路线。1488槽是E/\allowbreak{}G两目标、A/\allowbreak{}P两配置与3个seed的组合；A与P均为PyVRP 0.14.0的ILS，主要差异在罚项。G为带1.1米圆盘中心安全间距的有限合法图距离，取整到毫米，不是连续空间全局最短路，也不含航向、最小转弯半径或多车时空冲突。}
% content block 3:20 p
\bodytext{E/\allowbreak{}G矩阵在82/\allowbreak{}124田块相同；固定seed0的P\_\allowbreak{}E代表解在G上有正后悔13/\allowbreak{}124，42块矩阵不同者中29块仍无后悔。该诊断只针对固定点位与代表解，不能直接推断田块固有易错性；未来应将点位采样与并列解效应分开。}
% content block 3:8 h2
\Needspace{6\baselineskip}\subsection{2  已知失效与尚缺证据}
% content block 3:4 h3
\Needspace{4\baselineskip}\subsubsection{生成没有稳定胜出}
% content block 3:21 p
\bodytext{新20条装箱序列上，演化候选相对公开EoH为4胜、4平、12负，总计多用30箱。IID严格原始结果18/\allowbreak{}20，技术补齐后19/\allowbreak{}20；完整20例主比较仍为空。两种口径并列保留，不覆盖严格结果，不因缺失就宣布演化优于IID。[R4]}
% content block 3:5 h3
\Needspace{4\baselineskip}\subsubsection{选择收益存在明确边界}
% content block 3:22 p
\bodytext{BPP风险集是4组同前缀、不同后缀顺序的12条相关流，相对同一训练SBS平均多25.333箱。它给出部署前信息不足的明确反例。CVRP的24实例、432条历史记录中，0.25秒kNN改善约0.01441\%，仅由1例改选贡献；主CRC在小校准样本设置下结构性全回退。不能通过事后换风险阈值制造正结果。[R4]}
% content block 3:23 p
\bodytext{VRPTW 8特征选择器的描述性收益值得继续检查，但其选择时点尚未完全核明。如果8特征版本是在看到终测后才选定，它只能作为开发结果，必须用新冻结终测确认。没有母组差异和预定分析，不能直接说其显著或不显著。}
% content block 3:6 h3
\Needspace{4\baselineskip}\subsubsection{当前缺口对应的工作}
% content block 3:24 p
\bodytext{已有独立执行、恢复与重放入口，当前产品还需接线。后续技术部分统一规定应用、任务队列、地图与执行的接口；不以静态前端替代实时求解，也不重写可复用实验。}
% content block 3:7 h3
\Needspace{4\baselineskip}\subsubsection{待补的有限证据}
% content block 3:2 bullets
\begin{itemize}
\item 核对今天重跑的兼容性，补当前执行验收；历史PASS继续标为历史。
\item 建立跨包的母组、程序谱系和重复使用登记，避免重复计数。
\item 对缺失费用保留unknown；生成token、物理请求和全局成本不能凭汇总补齐。
\item 完整农业作业的联合验算、真实LLM在线编排与实际调用仍需贯通。
\end{itemize}
% content block 3:25 p
\bodytext{下一步围绕这些边界设计新的冻结确认；不重复把已有实验当作从零起步，也不将单批阴性外推到全部学习或生成方向。}
% content block 3:26 p
\bodytext{老师指南中的其他实践项目另行保留，与装箱生成的P1/\allowbreak{}P2批次编号不同：指南P1有120次汇总；指南P2合法排序360/\allowbreak{}360之外另有24次异常合同失败，不能写所有输入均正确；指南P4已有NAS与重训。它们不自动提供农业监督，也无需同时扩展成三条论文路线。[R2、R4]}
% content block 3:9 h2
\Needspace{6\baselineskip}\subsection{3  任务语义与算法能力合同}
% content block 3:27 p
\bodytext{统一输入、目标、能力声明和证据格式，不要求所有任务由一个求解器承担。未建模的关键约束必须拒绝；独立校验能发现错误，却不能补回编译时删去的可行域，也不能证明错误目标下的最优性。}
% content block 3:4 table
\begingroup
\fontsize{9.6}{14.6}\selectfont
\setlength{\tabcolsep}{4pt}
\renewcommand{\arraystretch}{1.18}
\begin{longtable}{@{}>{\raggedright\arraybackslash}p{\dimexpr 0.144092\linewidth-3.4582pt\relax}>{\raggedright\arraybackslash}p{\dimexpr 0.425072\linewidth-10.2017pt\relax}>{\raggedright\arraybackslash}p{\dimexpr 0.430836\linewidth-10.3401pt\relax}@{}}
\toprule
{\sffamily\bfseries 任务} & {\sffamily\bfseries 服务完成的判据} & {\sffamily\bfseries 适配与证据} \\
\midrule
\endfirsthead
\toprule
{\sffamily\bfseries 任务} & {\sffamily\bfseries 服务完成的判据} & {\sffamily\bfseries 适配与证据} \\
\midrule
\endhead
\midrule
\multicolumn{3}{r}{\fontsize{8}{10}\selectfont\color{Muted}续下页}\\
\endfoot
\bottomrule
\endlastfoot
点服务 & 在指定客户发生一次显式服务事件；经过其坐标不自动完成。 & 复用现有 E/\allowbreak{}G 流程及点服务 checker。 \\
\addlinespace[4pt]
整行作业 & 沿允许方向连续完成整条任务线，满足作业状态、通行与端点条件。 & 保存完整行服务路径，不用中点代替整行。 \\
\addlinespace[4pt]
面积覆盖 & 作业开启时，农具扫掠并集覆盖目标区域。 & 分别计算遗漏、重复和禁入侵入；Fields2Cover 可作为候选。 \\
\addlinespace[4pt]
连接与导航 & 任务间连接可通行，并满足声明的姿态、转弯、倒车或执行约束。 & 使用合法线段图、转弯连接或 OMPL；实机接口另设。 \\
\end{longtable}
\endgroup
% content block 3:28 p
\bodytext{每个行或区域服务模式绑定入口姿态、出口姿态、内部完整路径、成本和资源需求，不能分别挑最便宜的入口与出口后拼成不存在的服务。排序后须展开完整见证。转移若依赖进入弧、航向、时间、载荷、共享农具或历史碾压，标量 c(i,j) 可能不足，应扩展模式状态或使用整条路线评估器；静态容量和时间窗仍可与静态矩阵共存，只需正确累计资源和时间。}
% content block 3:29 p
\bodytext{算法登记表绑定版本、适配器、已测约束组合、车队语义、目标、精确性范围、输入输出与预算。未知能力按不支持处理。先筛能力，再在合格池中固定选择或学习推荐；无合格算法时明确缺口。车队上限不等于恰好 K 条非空路线，词典序目标不等于加权和。NL124 因 12 个单位需求、容量 5 和至多 3 车恰好迫使使用 3 车，不能据此宣称任意低负载场景已支持 exact-K。}
% content block 3:8 h3
\Needspace{4\baselineskip}\subsubsection{首个适配器的边界}
% content block 3:30 p
\bodytext{点服务 E 使用直线距离，G 使用有限合法线段图上的路径长度。已有 1.1 米圆盘中心安全间距检查不代表航向、转弯半径或多车时空避碰已实现。编辑几何后必须重建受影响的合法图与距离矩阵。随后分别接整行、区域覆盖和导航；Fields2Cover 的条带、地头与转弯流程不等于多车调度已解决，turn\_\allowbreak{}modes 的入弧语义也须先统一 depot、车数和服务定义。}
% content block 3:3 source
\sourcetext{工程来源 T13、T17、T18、T28，见技术来源附录。}
% content block 3:10 h2
\Needspace{6\baselineskip}\subsection{4  农业主合同与四类案例}
% content block 3:31 p
\bodytext{建议以现有点服务或果园任务贯通真实执行，将首个深入农业实验限定为既定作业行、明确连接模式、单车和固定站补给。它可与学习型路线改进、自动设计或预算决策结合；一次实验只改变一个核心机制。四类案例分别验收，不要求全部进入同一篇论文。}
% content block 3:9 h3
\Needspace{4\baselineskip}\subsubsection{深入合同的最小边界}
% content block 3:3 bullets
\begin{itemize}
\item 输入为有来源的米制田块、既定作业行、车辆足迹与转弯参数、每行耗料和固定站。
\item 每行不可拆，必须完整服务；允许方向由任务定义，服务与非作业移动明确区分。
\item 只研究一种资源，回站补满，补给时间与行驶时间同时累计；超容量行显式判不支持或不可行。
\item 目标至少分解为作业、空驶、转弯、回仓和补给成本；优化主指标预先冻结。
\end{itemize}
% content block 3:10 h3
\Needspace{4\baselineskip}\subsubsection{课题案例各自验证什么}
% content block 3:5 table
\begingroup
\fontsize{9.6}{14.6}\selectfont
\setlength{\tabcolsep}{4pt}
\renewcommand{\arraystretch}{1.18}
\begin{longtable}{@{}>{\raggedright\arraybackslash}p{\dimexpr 0.165706\linewidth-3.9769pt\relax}>{\raggedright\arraybackslash}p{\dimexpr 0.443804\linewidth-10.6513pt\relax}>{\raggedright\arraybackslash}p{\dimexpr 0.390490\linewidth-9.3718pt\relax}@{}}
\toprule
{\sffamily\bfseries 农业案例} & {\sffamily\bfseries 功能验收重点} & {\sffamily\bfseries 论文角色} \\
\midrule
\endfirsthead
\toprule
{\sffamily\bfseries 农业案例} & {\sffamily\bfseries 功能验收重点} & {\sffamily\bfseries 论文角色} \\
\midrule
\endhead
\midrule
\multicolumn{3}{r}{\fontsize{8}{10}\selectfont\color{Muted}续下页}\\
\endfoot
\bottomrule
\endlastfoot
果园运输 & 合法导航、载荷任务与路径回放，推荐到真实运行贯通。 & 现有组件复用与入口回归。 \\
\addlinespace[4pt]
农田巡检避障 & 足迹、障碍与局部决策合法；动态场景采用固定可复现时序。 & 课题能力；是否进入主文取决于机制关联。 \\
\addlinespace[4pt]
多作业点规划 & 容量/\allowbreak{}车队/\allowbreak{}时间规则和强基线；点任务完整服务。 & 已有资产接续及合同外检查。 \\
\addlinespace[4pt]
不规则地块覆盖 & 工具扫掠、漏作、重叠、转弯与完整工作量。 & 课题能力；提供实质新机制时再纳入论文主实验。 \\
\end{longtable}
\endgroup
% content block 3:32 p
\bodytext{真实田块边界不自动带来真实行宽、需求、载荷、地头或障碍。逐字段标为实测、原件重建、文献设定、合成或未知。真实来源几何上的合成作业应称半合成，不借其他论文的田间验证为本项目背书。}
% content block 3:33 p
\bodytext{点服务、行服务与面积覆盖可以共享任务身份、注册、运行和证据格式，但不能无损压成同一个标量距离矩阵。转弯模式、入出姿态、时间或资源可能改变后续可行性，必须由模式状态或整条路线评估器保留。}
% content block 3:4 source
\sourcetext{不强制四类案例共用无损距离矩阵，也不把四类分数直接平均成“农业能力总分”。}
% content block 3:11 h2
\Needspace{6\baselineskip}\subsection{5  服务与运动语义的资格实验}
% content block 3:34 p
\bodytext{农业增量首先来自正确的作业语义。已有Fields2Cover、HGS-CARP/\allowbreak{}NEARP、PyVRP reload与turn\_\allowbreak{}modes提供可用底座，尚缺本项目完整行作业、补给和扫掠联合验收。下列反例先用小规模精确或穷举参照通过，再生成训练标签与论文结果。}
% content block 3:6 table
\begingroup
\fontsize{9.6}{14.6}\selectfont
\setlength{\tabcolsep}{4pt}
\renewcommand{\arraystretch}{1.18}
\begin{longtable}{@{}>{\raggedright\arraybackslash}p{\dimexpr 0.187320\linewidth-4.4957pt\relax}>{\raggedright\arraybackslash}p{\dimexpr 0.403458\linewidth-9.6830pt\relax}>{\raggedright\arraybackslash}p{\dimexpr 0.409222\linewidth-9.8213pt\relax}@{}}
\toprule
{\sffamily\bfseries 语义区别} & {\sffamily\bfseries 最小资格例} & {\sffamily\bfseries 保存证据} \\
\midrule
\endfirsthead
\toprule
{\sffamily\bfseries 语义区别} & {\sffamily\bfseries 最小资格例} & {\sffamily\bfseries 保存证据} \\
\midrule
\endhead
\midrule
\multicolumn{3}{r}{\fontsize{8}{10}\selectfont\color{Muted}续下页}\\
\endfoot
\bottomrule
\endlastfoot
至多K与恰K & 2个单位需求、Q=2、K=2；一条路线仅满足至多K，两条单服务路线满足恰K。 & 可行域见证与独立车数检查；NL124不能识别此差异。 \\
\addlinespace[4pt]
经过与服务 & 路径经过客户但无服务动作。 & 显式服务账本、每任务一次服务及合法替代。 \\
\addlinespace[4pt]
行端与整行 & 两端可达，但中段未开机服务。 & 行区间、方向、作业开关和整行见证。 \\
\addlinespace[4pt]
中心线与整机 & 中心线合法，但车体或偏置农具越界。 & 车体占用、农具包络和作业扫掠分别检查。 \\
\addlinespace[4pt]
补给与时间 & 回站补满，但总时间不得清零；单行耗料超过容量。 & 资源和事件时间账，超容量明确拒绝或判不可行。 \\
\addlinespace[4pt]
单位与任务变化 & ft转m保持物理任务；改变幅宽可改变可行域。 & 原单位、换算来源与用户确认，不混作同一种迁移。 \\
\end{longtable}
\endgroup
% content block 3:11 h3
\Needspace{4\baselineskip}\subsubsection{固定行的首个最小实验}
% content block 3:35 p
\bodytext{先固定一车一趟、刚性机具、给定作业行和已验连接目录；地头只供关闭机具后的连接。候选仅提议翻向、搬移、交换或有界重排，底层连接与checker固定。对照Fields2Cover端到端方案、共用真实连接成本的强传统序列搜索、自动配置与多启动；分别报空驶/\allowbreak{}转弯和包含固定行作业的总成本。}
% content block 3:36 p
\bodytext{已有JPDS-NN研究作业行与入口联合学习，连续曲率地头研究也已处理整机与机具几何。行序学习与转弯连接均有直接近邻；本项目应比较真实连接和资源合同下的增量，而不把这些模块本身视为空白。\href{https://arxiv.org/html/2503.02369v1}{JPDS-NN}；\href{https://onlinelibrary.wiley.com/doi/10.1002/rob.22489}{连续曲率地头覆盖}。}
% content block 3:37 p
\bodytext{模型逐步增加单站补给后，保存服务、行驶、转弯、回仓、补给的事件级账。PyVRP源码已有reload\_\allowbreak{}depots和max\_\allowbreak{}reloads，现URP适配器设置为0，需显式启用与适配；P102历史工具单次使用、回仓换具不等于补货重用或含服务时间的补给闭环。}
% content block 3:5 source
\sourcetext{农业合同与覆盖来源：依据：\href{https://github.com/Fields2Cover/Fields2Cover/blob/93c0989ab2fc4a3f3d0a18bbd224d0e267a03fcc/src/fields2cover/route\_planning/route\_planner\_base.cpp\#L159-L190}{Fields2Cover路线入口}；\href{https://github.com/Fields2Cover/Fields2Cover/blob/884d895b59192882476e986ba44ea9143a06a6a9/include/fields2cover/types/GenericImplement.h\#L20-L56}{v3机具}；\href{https://github.com/Fields2Cover/Fields2Cover/blob/93c0989ab2fc4a3f3d0a18bbd224d0e267a03fcc/tests/python/path\_planning/PP\_checker.py\#L19-L26}{Python检查}与\href{https://github.com/Fields2Cover/Fields2Cover/blob/93c0989ab2fc4a3f3d0a18bbd224d0e267a03fcc/tests/cpp/test\_helpers/path\_planning\_checker.hpp\#L71-L93}{C++检查}；\href{https://zenodo.org/records/14524735}{田界数据}。小例为本项目合同推导，未运行实验。}
% content block 3:12 h2
\Needspace{6\baselineskip}\subsection{6  转向与跨合同适应的可执行入口}
% content block 3:38 p
\bodytext{turn\_\allowbreak{}modes已有入弧状态最短路、start/\allowbreak{}end/\allowbreak{}conn预处理、fixed\_\allowbreak{}sequence模式优化和evaluate(instance,candidate)入口。36例来自12根图×3规则，与精确oracle一致；丢方向对照32例错误。它精确求解的是给定服务顺序的模式与路径，不是整体任务分配和顺序最优；已有services事件仍是节点服务。}
% content block 3:39 p
\bodytext{现入口每条路线出发无前序入弧，回仓后结束，不计末入弧与下一次出发的闭合转向；depot\_\allowbreak{}reset未成为可切换且强制校验的通用合同。subset\_\allowbreak{}dp没有车辆数维度，exact-K、行服务与覆盖不能由NL124或字段名称自动继承。HGS-CARP已有首末服务模式的子序列成本结构，复用它是经典建模，不当作新方法。}
% content block 3:12 h3
\Needspace{4\baselineskip}\subsubsection{最小实验与强对照}
% content block 3:40 p
\bodytext{合同适配实验固定正确模式解释器，仅开放一个车数、depot语义或成本接口，比较手写适配、冻结检索、标量反馈与结构化反例。目标是合同通过、旧能力回归与修复总成本。}
% content block 3:41 p
\bodytext{启发式迁移实验保持任务服务与车队/\allowbreak{}depot语义，改变有作业依据的转弯成本或禁转；固定DP和邻域，仅开放优先级时属于搜索控制。强对照是完全不改代码、直接调用新evaluate的通用mode-aware搜索，另比目标从零、冻结检索、一次手工适配与同空间自动配置。}
% content block 3:13 h3
\Needspace{4\baselineskip}\subsubsection{数据条件和去留}
% content block 3:42 p
\bodytext{同根图换规则用于隔离受控变化，未见根图与机具参数组合用于外推；12个现有根图仅作开发和资格。半径、地头宽度与工具参数应来自真实作业设定。若只换正确cost就完成适配，或强通用搜索已等效，保留工程收益并收缩知识迁移主张；总成本改善且未见任务合法质量维持时，再扩大合同与规模。}
% content block 3:6 source
\sourcetext{转向实现依据：本地：urp\_\allowbreak{}research/\allowbreak{}turn\_\allowbreak{}modes/\allowbreak{}solver.py:13、evaluate\_\allowbreak{}candidate.py:10；HGS-CARP/\allowbreak{}Program/\allowbreak{}SeqData.h:58、Individu.cpp:320。公开对照：\href{https://github.com/vidalt/HGS-CARP/blob/08fde6a5672c3814dca1b4450aa9e5968b1d273a/Program/Params.cpp\#L1041-L1097}{HGS成本与turn}；\href{https://github.com/vidalt/HGS-CARP/blob/08fde6a5672c3814dca1b4450aa9e5968b1d273a/Program/SeqData.cpp\#L514-L538}{子序列mode表}。}
% content block 4:4 h1
\clearpage
\hypertarget{part_3}{}
\section{三 地图 系统与执行实施}
% content block 4:13 h2
\Needspace{6\baselineskip}\subsection{1  技术路线与现有工程接入}
% content block 4:43 p
\bodytext{实施路线采用 OpenLayers 的稳定二维地图编辑、pyproj 与 Shapely 的后端坐标和几何层，保留 React/\allowbreak{}Vite 与 SQLite。地图、研究 CLI、算法适配器和独立校验器共享任务合同与版本身份；先接通可编辑输入、真实执行、独立重算和证据回放，再按实测需求扩展覆盖、影像和三维。以下接口与验收均为拟定实施要求。}
% content block 4:14 h3
\Needspace{4\baselineskip}\subsubsection{先复用已有资产}
% content block 4:44 p
\bodytext{截至技术方案所载的 2026年10月11日只读核查，当前 React 仍读取静态 JSON，HTTP 后端声明 executes\_\allowbreak{}experiments=false；数据库有 5 个实验 ID、28 个 artifact，接口以导入、查询和回放为主。CVRP 核心与独立 checker 已迁入，但尚未接入当前界面；未找到统一派发、取消与恢复接口。此状态不否定独立研究实验已经执行，也不代表保存结果已在本轮重跑。}
% content block 4:45 p
\bodytext{存档 planner.html 和 planner.js 已有点地图编辑、坐标与需求增删、浏览器基线计算、导入导出和文件桥接。迁移应保留可用操作、样例与回归测试，再接入当前 React 状态管理。独立 CLI、持久槽位和个别实验恢复驱动一并复用；存档原型、当前应用、可运行环境和保存证据分别验收。}
% content block 4:15 h3
\Needspace{4\baselineskip}\subsubsection{模块责任}
% content block 4:4 bullets
\begin{itemize}
\item React 管理表单、面板和应用状态；OpenLayers 实例由独立适配层管理。临时编辑对象与已提交 SceneSpec 分离，使用稳定对象 ID 和明确提交事务，避免渲染重建地图或可变几何改写冻结输入。
\item 后端完成坐标、拓扑、业务校验及场景编译；编译输出图、矩阵、服务模式或覆盖实例，并声明精确等价、限制或近似。规划器负责寻找方案，独立 checker 重算约束与成本。
\item runner 使用受监管子进程调用已有 CLI 或适配器，控制预算、取消和输出。首版单机单 worker，保留现有导入、查询和回放接口；不能只把 executes\_\allowbreak{}experiments 改为 true。
\end{itemize}
% content block 4:46 p
\bodytext{实施前先恢复一个可重复运行的 CLI 最小案例，核对解释器、Python 和原生 GIS/\allowbreak{}规划依赖、数据路径、SQLite 实际链接版本及发行版安全补丁。历史结果可读取与今天的运行链可用，是两项验收。}
% content block 4:7 source
\sourcetext{工程来源 T14、T16、T19、T21、T28，见技术来源附录。}
% content block 4:14 h2
\Needspace{6\baselineskip}\subsection{2  地图引擎选型与替代条件}
% content block 4:47 p
\bodytext{主工作台选用 OpenLayers：现成 Draw、Modify、Snap 适合真实地理底图、米制局部平面和几何编辑。工具栏、撤销和业务校验仍由应用实现；首版核心编辑避开实验性 WebGL API。地图引擎只负责显示与交互，不自动提供容量、禁转、作业宽度或碰撞约束。}
% content block 4:7 table
\begingroup
\fontsize{9.6}{14.6}\selectfont
\setlength{\tabcolsep}{4pt}
\renewcommand{\arraystretch}{1.18}
\begin{longtable}{@{}>{\raggedright\arraybackslash}p{\dimexpr 0.187320\linewidth-4.4957pt\relax}>{\raggedright\arraybackslash}p{\dimexpr 0.389049\linewidth-9.3372pt\relax}>{\raggedright\arraybackslash}p{\dimexpr 0.423631\linewidth-10.1671pt\relax}@{}}
\toprule
{\sffamily\bfseries 方案} & {\sffamily\bfseries 适用条件} & {\sffamily\bfseries 本项目取舍} \\
\midrule
\endfirsthead
\toprule
{\sffamily\bfseries 方案} & {\sffamily\bfseries 适用条件} & {\sffamily\bfseries 本项目取舍} \\
\midrule
\endhead
\midrule
\multicolumn{3}{r}{\fontsize{8}{10}\selectfont\color{Muted}续下页}\\
\endfoot
\bottomrule
\endlastfoot
OpenLayers & 真实底图、专业投影、局部 XY 和栅格需要统一。 & 首选稳定二维编辑；锁版本后做触控、撤销和实例销毁重建测试。 \\
\addlinespace[4pt]
MapLibre 加 Terra Draw & 界面主要面向经纬度展示与 WebGL 矢量底图。 & 窄原型备选；先验证投影、编辑、触控与资源成本。常用 COG 插件仅支持 EPSG:3857。 \\
\addlinespace[4pt]
Geoman 或 Leaflet & 简单采集页，或所需编辑操作能由现成插件覆盖。 & Geoman 按选定引擎核对 Free/\allowbreak{}Pro 功能及许可；Leaflet 不作为复杂主工作台首选。 \\
\addlinespace[4pt]
deck.gl & 稳定分层、按需加载和工作集限制后，大数据渲染仍不达标。 & 增加只读高密度图层；不把社区 editable-layers 当核心编辑依赖。 \\
\addlinespace[4pt]
CesiumJS & 真实高程、坡地或三维问题确实影响决策。 & 条件性扩展；二维拓扑编辑仍需自建。 \\
\end{longtable}
\endgroup
% content block 4:48 p
\bodytext{Terra Draw 的 README 兼容表与适配器包依赖范围存在差异，不能仅凭表格断言 MapLibre 6 不支持。锁定地图核心、编辑核心、适配器和工具栏的具体版本，实测绘制、撤销、触控及销毁重建，再决定替代。}
% content block 4:49 p
\bodytext{MapStore 适合日后的多地图、权限和插件平台，当前引入会扩大重构。QGIS/\allowbreak{}QField 作为专业配准与离线采集外部工具接文件；Rerun 地图当时只支持点和线、不支持面，Foxglove 自 2024 年停止 Studio 开源版开发，不作为主编辑器。}
% content block 4:16 h3
\Needspace{4\baselineskip}\subsubsection{选型试验如何结束}
% content block 4:50 p
\bodytext{已有资产为存档编辑器及当前 React 页面。最小实验是在同一真实田块与固定 seed 合成集上完成导入、顶点修改、撤销、保存重载和路线回放；强简单基线为 OpenLayers 稳定二维层与数值表单。仅在同一脚本、设备、正确性检查和数据量下，MapLibre 加 Terra Draw 在手机操作、内存或维护成本上出现可重复优势时才替换主选。不能仅用官网演示或不同数据的帧率作决定。}
% content block 4:8 source
\sourcetext{工程来源 T1、T2、T3、T4、T5、T6、T7、T15，见技术来源附录。}
% content block 4:15 h2
\Needspace{6\baselineskip}\subsection{3  坐标精度与安全几何}
% content block 4:51 p
\bodytext{每个场景必须声明唯一计算 CRS 和单位。真实田块采用适合当地的米制投影或经测量定义的局部坐标；合成场景采用 local\_\allowbreak{}xy 米制坐标。权威计算与持久化保留完整 CRS、原始输入和精度；专业视图可使用工作坐标，把底图重投影显示。}
% content block 4:8 table
\begingroup
\fontsize{9.6}{14.6}\selectfont
\setlength{\tabcolsep}{4pt}
\renewcommand{\arraystretch}{1.18}
\begin{longtable}{@{}>{\raggedright\arraybackslash}p{\dimexpr 0.216138\linewidth-5.1873pt\relax}>{\raggedright\arraybackslash}p{\dimexpr 0.353026\linewidth-8.4726pt\relax}>{\raggedright\arraybackslash}p{\dimexpr 0.430836\linewidth-10.3401pt\relax}@{}}
\toprule
{\sffamily\bfseries 表示} & {\sffamily\bfseries 用途} & {\sffamily\bfseries 不可混淆的边界} \\
\midrule
\endfirsthead
\toprule
{\sffamily\bfseries 表示} & {\sffamily\bfseries 用途} & {\sffamily\bfseries 不可混淆的边界} \\
\midrule
\endhead
\midrule
\multicolumn{3}{r}{\fontsize{8}{10}\selectfont\color{Muted}续下页}\\
\endfoot
\bottomrule
\endlastfoot
米制 XY 与完整 CRS & 距离、缓冲、碰撞、路线成本和安全裕量。 & 米制不等于无投影尺度畸变，也不等于影像或定位有厘米精度。 \\
\addlinespace[4pt]
RFC 7946 GeoJSON & WGS84 经度、纬度顺序的交换或地理展示。 & 不得将米制 XY 冒充标准 GeoJSON；逐来源核对轴序。 \\
\addlinespace[4pt]
视图坐标与像素 & 绘制、捕捉与展示简化。 & 屏幕捕捉容差不等于物理安全距离；简化副本不写回权威几何。 \\
\end{longtable}
\endgroup
% content block 4:52 p
\bodytext{源方案所核 OpenLayers v10.11.0 文档提示内建 UTM 转换误差量级约 0.1 米。精确编辑需注册并验证实际采用的 proj4js 转换，用参考控制点与后端 PROJ 对照；往返转换一致并不能证明绝对位置正确。经纬度六位小数也不是 1 厘米：纬度格距约 0.111 米，四舍五入单轴误差可达约 0.056 米，经度误差随纬度变化。}
% content block 4:53 p
\bodytext{pyproj 应明确轴顺序、变换适用范围与精度状态，离线包携带实际需要的 PROJ 数据和格网。Web Mercator 的显示长度不能直接用于农具宽度或安全距离；Shapely 的 XY 判断不会因附加 Z 自动成为三维碰撞检查。}
% content block 4:17 h3
\Needspace{4\baselineskip}\subsubsection{合法性以后端结果为准}
% content block 4:54 p
\bodytext{输入检查自交、空几何、重复点、洞的位置、最小通道和任务点是否处于可用区域。Shapely make\_\allowbreak{}valid 可能产生多面或几何集合，修复应展示差异并经明确接受，不能静默改变任务边界。拖动时可做快速预览，提交后以后端校验确定能否运行。}
% content block 4:55 p
\bodytext{安全区域由车体与农具包络、安全裕量及误差预算推导。圆盘可作为首版近似，但带航向的长方形车体、转弯半径、坡度及多车时空碰撞需要各自模型。验收误差阈值由数据精度和安全裕量决定，不凭地图缩放等级或动画观感设定。}
% content block 4:56 p
\bodytext{坐标资格检查先使用有参考控制点的小区域和固定合成几何，以后端 pyproj/\allowbreak{}Shapely 与数值输入作为简单基线。只有绝对位置、长度和安全裕量误差均在声明范围内，才开放精细编辑；没有坐标来源或校准证据的数据保留未知状态，不猜 EPSG。}
% content block 4:9 source
\sourcetext{工程来源 T2、T3、T8、T9，见技术来源附录。}
% content block 4:16 h2
\Needspace{6\baselineskip}\subsection{4  真实地图与合成输入}
% content block 4:57 p
\bodytext{底图是背景，经过确认的工作边界和任务对象才进入规划。Fields2Benchmark 的 350 块田界具有真实来源，但公开说明不足以确认所有 WKT 的统一 CRS 和现场精度，立陶宛来源还限非商业。导入分别记录真实来源、是否配准与精度状态，不把公开边界直接当精确实测作业地图。}
% content block 4:9 table
\begingroup
\fontsize{9.6}{14.6}\selectfont
\setlength{\tabcolsep}{4pt}
\renewcommand{\arraystretch}{1.18}
\begin{longtable}{@{}>{\raggedright\arraybackslash}p{\dimexpr 0.208934\linewidth-5.0144pt\relax}>{\raggedright\arraybackslash}p{\dimexpr 0.453890\linewidth-10.8934pt\relax}>{\raggedright\arraybackslash}p{\dimexpr 0.337176\linewidth-8.0922pt\relax}@{}}
\toprule
{\sffamily\bfseries 输入} & {\sffamily\bfseries 处理要求} & {\sffamily\bfseries 保留证据} \\
\midrule
\endfirsthead
\toprule
{\sffamily\bfseries 输入} & {\sffamily\bfseries 处理要求} & {\sffamily\bfseries 保留证据} \\
\midrule
\endhead
\midrule
\multicolumn{3}{r}{\fontsize{8}{10}\selectfont\color{Muted}续下页}\\
\endfoot
\bottomrule
\endlastfoot
GeoJSON 与 GeoPackage & GeoJSON 按 WGS84 读取；完整有 CRS 的矢量交付优先 GeoPackage。 & 原文件、源 CRS、图层和属性映射。 \\
\addlinespace[4pt]
Shapefile 与 CSV & Shapefile 成套打包并核 .prj；CSV 明确 XY/\allowbreak{}经纬度、轴序、单位。 & 缺失字段、任何人工映射及导入报告。 \\
\addlinespace[4pt]
GeoTIFF 与 COG & 小图可直接显示，大图按切片或 COG 分块和概览读取。 & CRS、地理变换、分辨率、NoData、采集时间。 \\
\addlinespace[4pt]
普通图片 & 无地理参考时先做控制点配准；仅有已知尺寸时作为合成平面参考。 & 控制点、模型、残差及适用范围。 \\
\addlinespace[4pt]
合成场景 & 固定生成器版本、seed、范围、障碍与任务分布；编辑另存修订。 & 原始 seed、生成参数、编辑后几何。 \\
\end{longtable}
\endgroup
% content block 4:58 p
\bodytext{COG 是方便按范围读取的栅格组织方式，不保证无限影像都能在浏览器中加载。远端应支持 HTTP Range 字节读取和所需 CORS；颜色、重投影、缓存与读取范围须在目标手机实测。不满足时使用本地瓦片服务或预处理的有限区域包。}
% content block 4:59 p
\bodytext{离线只打包有授权的数据。PMTiles 或后端离线包须连同样式字体、精灵图及其他资源处理；首版限定选定项目区域和容量。OpenStreetMap 官方栅格与矢量瓦片服务禁止为离线批量预取，开源地图库不代表底图可复制。保留来源、署名、许可用途、下载时间与包校验值。}
% content block 4:60 p
\bodytext{显示与分析分开：类别图用最近邻重采样，DEM 保留原始高程和高程基准，NoData 视为未知。显示归一化、RGB 着色和降分辨率不替代分析像元。公开示例若含自定义类型或非标准轴序，须按来源单独转换；树点不能擅自连成“实测作业行”，无田界数据不能自动补成实测边界。}
% content block 4:61 p
\bodytext{大图资格试验从有限区域、原始像元和本地瓦片简单基线开始，逐级增加文件大小与可见顶点。只有目标手机的内存、首次呈现、按需读取及离线恢复均达标才扩大工作集；处理和显示副本必须可从原件重建。}
% content block 4:10 source
\sourcetext{工程来源 T2、T9、T10、T11、T12、Tinteraction2，见技术来源附录。}
% content block 4:17 h2
\Needspace{6\baselineskip}\subsection{5  桌面手机交互与部署边界}
% content block 4:62 p
\bodytext{Ubuntu 首版采用响应式 Web 应用，本机保存原始数据；手机通过经认证的局域网入口访问，默认不公开执行端口。桌面与手机共享场景格式和状态机，布局分别适配。PWA 如引入，只作为安装与缓存增强，其离线范围应逐项验收；不能据此宣称 Python 求解器、原生 GIS 依赖或完整实验已能在手机独立运行。}
% content block 4:63 p
\bodytext{手机完全脱离 Ubuntu 是单独交付目标，需要明确任务包、可用求解能力、存储配额、数据导出与恢复流程；浏览器缓存不能作为唯一原件。桌面薄壳放在闭环之后：Tauri 2 须测 Ubuntu WebKitGTK 与地图渲染，Electron 内置 Chromium 可作兼容备选；两者都需要单独打包与监管 Python 进程。}
% content block 4:18 h3
\Needspace{4\baselineskip}\subsubsection{完整操作路径}
% content block 4:5 bullets
\begin{itemize}
\item 导入或选择场景，确认 CRS、单位、比例尺和底图来源，显示数据质量与边界检查结果。
\item 编辑田块、洞、障碍、客户、仓库和车辆；浏览、绘制、修改、删除模式明确分开，属性可数值输入。
\item 按能力筛选规划器，展示任务类型、参数、seed、预算和限制；提交前保存新修订，错误定位到对象。
\item 运行期间允许浏览或另开修订；完成后按对应输入版本回放，地图变更不得自动套用旧解。
\end{itemize}
% content block 4:64 p
\bodytext{基础编辑应包含选择、顶点拖动、增删顶点、撤销重做、吸附开关、数值输入和取消；吸附显示目标对象，移动保留稳定 ID。手机使用明确“完成”“取消”按钮，不依赖双击、右键或修饰键；地图拖动与顶点拖动隔离，提供放大定位，常用触控目标拟至少 44 CSS 像素。删除与离开未保存修订须可恢复。复杂切割、合并和共边编辑后置。}
% content block 4:19 h3
\Needspace{4\baselineskip}\subsubsection{结果解释与可追溯修改}
% content block 4:65 p
\bodytext{分层显示底图、原始边界、可通行区域、任务、候选路线和已核验路线。车辆用颜色加线型或编号区分；选中路段显示距离、服务顺序、载荷、代价来源和校验状态，支持逐车开关和违规点定位。失败车辆、未分配服务与未检查约束不能被成功路线遮住。}
% content block 4:66 p
\bodytext{访问顺序、几何路径与带时间轨迹分开；无时间戳时只按路径进度播放，不称真实轨迹，不据此证明多车无冲突。算法比较固定场景哈希、目标、车队语义、预算和视口。地图、成本图表与谱系通过 scene\_\allowbreak{}id、修订、input\_\allowbreak{}hash、run\_\allowbreak{}id、candidate/\allowbreak{}source\_\allowbreak{}hash 联动，指标及失败分母由后端统一。}
% content block 4:67 p
\bodytext{拖动路线默认生成途经点、禁区、走廊偏好或锁定顺序等新约束后重算；自由手绘先标未验证，通过同等级独立校验后才能提升状态。每次编辑记录原意图、旧新任务版本、补丁、确认来源与 checker 结果，并区分纠正理解、改变任务、更新偏好和表达不确定；接受、拖动、撤销不自动成为算法优劣或最优标签。}
% content block 4:68 p
\bodytext{可视编辑的最小实施试验先固定 solver 和有错任务样例，比较地图编辑与字段表单、固定澄清和规则修改器；以同一任务版本与回放日志检查错误是否定位和消除。离线编辑回放只能证明接口有效；减少用户漏错和节省时间还需真实参与者资料，研究试验另列。字段表单已满足任务时不因展示效果增加模型。}
% content block 4:11 source
\sourcetext{工程来源 T12、T14、T26、Tinteraction1，见技术来源附录。}
% content block 4:18 h2
\Needspace{6\baselineskip}\subsection{6  数据接口 执行状态与安全}
% content block 4:69 p
\bodytext{以下为新增接口的最低合同，并非当前 API 已实现的清单。先冻结请求样例、响应结构、错误码和身份字段，再由前后端共同实现；GUI 与批量 CLI 使用同一冻结输入、runner 和 checker。}
% content block 4:10 table
\begingroup
\fontsize{9.6}{14.6}\selectfont
\setlength{\tabcolsep}{4pt}
\renewcommand{\arraystretch}{1.18}
\begin{longtable}{@{}>{\raggedright\arraybackslash}p{\dimexpr 0.259366\linewidth-4.1499pt\relax}>{\raggedright\arraybackslash}p{\dimexpr 0.740634\linewidth-11.8501pt\relax}@{}}
\toprule
{\sffamily\bfseries 对象或接口} & {\sffamily\bfseries 最低内容与责任} \\
\midrule
\endfirsthead
\toprule
{\sffamily\bfseries 对象或接口} & {\sffamily\bfseries 最低内容与责任} \\
\midrule
\endhead
\midrule
\multicolumn{2}{r}{\fontsize{8}{10}\selectfont\color{Muted}续下页}\\
\endfoot
\bottomrule
\endlastfoot
SceneSpec & 格式版本、scene\_\allowbreak{}id、scene\_\allowbreak{}revision、基准修订、CRS 与单位；几何、任务、车辆、目标顺序、整数缩放；input\_\allowbreak{}hash 与来源。 \\
\addlinespace[4pt]
POST /\allowbreak{}scenes/\allowbreak{}validate & 坐标、拓扑、业务校验；返回对象 ID、错误位置、物理裕量及可运行状态。 \\
\addlinespace[4pt]
POST /\allowbreak{}jobs & 冻结 scene\_\allowbreak{}revision 和 input\_\allowbreak{}hash；接收 planner\_\allowbreak{}id、参数、seed、预算与幂等键，返回 job\_\allowbreak{}id。 \\
\addlinespace[4pt]
任务状态与取消 & 绑定 job、attempt、worker、输入版本及请求序号；区分求解状态和校验状态。 \\
\addlinespace[4pt]
ResultBundle & 算法与环境版本、输入哈希、原始输出、规范化路线、独立校验、指标、附件校验值。 \\
\addlinespace[4pt]
推荐和特征记录 & FeatureRecord 记特征时点、缺失与成本；RecommendationRecord 记算法 ID、参数、模型版本、选择和回退依据。 \\
\end{longtable}
\endgroup
% content block 4:20 h3
\Needspace{4\baselineskip}\subsubsection{可中断且可恢复的状态机}
% content block 4:70 p
\bodytext{持久任务状态至少包含排队、运行、请求取消、已取消、成功、失败、超时和中断。查询服务与执行环境分离，重计算不放在 HTTP 请求或普通后台回调。FastAPI/\allowbreak{}Pydantic 可用于接口和合同校验，但不能替代任务队列、进程监管或安全隔离。}
% content block 4:6 bullets
\begin{itemize}
\item 领任务采用短事务，记录所有权和到期处理；失效 worker 的迟到输出不能覆盖当前结果。重复提交用幂等键去重，旧场景结果只入历史。
\item 取消先持久记录请求，再通知算法或终止受监管进程组，确认退出后才能记已取消；页面停止等待不等于求解器已停止。
\item 崩溃恢复先核对进程、状态与结果包，再决定继续、导入或新建 attempt；重试不覆盖旧日志，也不默认追加模型调用预算。
\item 结果先写临时目录，经结构检查后原子封存和登记。非法候选也留存，独立校验通过后才发布。SQLite 单文件只有一个写者，运行中数据库不放到网络共享盘供多机同时写。
\end{itemize}
% content block 4:21 h3
\Needspace{4\baselineskip}\subsubsection{生成 评测和在线推荐分开}
% content block 4:71 p
\bodytext{生成请求只能访问已授权模型接口，不能顺手执行回答代码。候选在无密钥、默认断网、只读输入、专用输出目录的独立环境中运行，限制 CPU、内存、进程数、时间和磁盘；可信 checker 在候选之外重算。语法白名单只是准入，容器、Python 子进程或 Ray 资源声明也不自动构成完整隔离，需核验权限、挂载、网络、进程组清理和限制确实生效。}
% content block 4:72 p
\bodytext{在线只选择已登记且通过准入的算法；自然语言修改只生成结构化场景草案、缺失字段和歧义项，经几何与业务检查后形成新修订，不编造车宽、能耗或任务事实。生成器不得修改测试集、checker、目标或正式发布指针，候选不能自我入库。推荐只使用决策时已知信息，构图、特征及试跑计入在线成本，对不确定或分布外输入采用已验算的固定回退。}
% content block 4:22 h3
\Needspace{4\baselineskip}\subsubsection{证据和预算保留完整}
% content block 4:73 p
\bodytext{分别记录排队、几何预处理、构图、模型生成、求解、校验和端到端墙钟时间，保留失败、超时、非法、未评测和重试分母；未评测不是零质量。固定迭代数与相同墙钟预算分开，历史 250 次迭代不改称公平墙钟比较。}
% content block 4:74 p
\bodytext{JSONL 保存原始事件，文件包保存冻结输入、源码、路线见证和日志，SQLite 保存状态和元数据，Parquet 可作为分析派生表。生成记录关联父代、提示、原始响应、代码、评分及失败，保留全部程序身份。追踪需求扩大后可加入 MLflow，但它不替代语义合同与原始证据。}
% content block 4:75 p
\bodytext{本地与云端共用输入、结果 manifest 和收据；回收核对文件哈希、记录数、导入状态和独立检查报告，平台 success 只证明执行状态。原始包确认回收并满足留存要求后再清理远端暂存。}
% content block 4:12 source
\sourcetext{工程来源 T14、T19、T20、T21，见技术来源附录。}
% content block 4:19 h2
\Needspace{6\baselineskip}\subsection{7  地图与执行闭环验收}
% content block 4:76 p
\bodytext{首轮沿用 NL124 点服务协议完成“改输入、真运行、独立重算、回放”：每例 12 个单位需求、单车容量 5、车队上限 3。此例必然用满 3 车，只验收该合同，不扩称任意 exact-K、面积覆盖或实机协同。沿既有分工，前端迁移与交互由 GrokBot 推进，后端合同、执行与校验由 Hill 对接，先共同冻结接口样例和异常约定。}
% content block 4:11 table
\begingroup
\fontsize{9.6}{14.6}\selectfont
\setlength{\tabcolsep}{4pt}
\renewcommand{\arraystretch}{1.18}
\begin{longtable}{@{}>{\raggedright\arraybackslash}p{\dimexpr 0.194524\linewidth-3.1124pt\relax}>{\raggedright\arraybackslash}p{\dimexpr 0.805476\linewidth-12.8876pt\relax}@{}}
\toprule
{\sffamily\bfseries 验收阶段} & {\sffamily\bfseries 通过证据} \\
\midrule
\endfirsthead
\toprule
{\sffamily\bfseries 验收阶段} & {\sffamily\bfseries 通过证据} \\
\midrule
\endhead
\midrule
\multicolumn{2}{r}{\fontsize{8}{10}\selectfont\color{Muted}续下页}\\
\endfoot
\bottomrule
\endlastfoot
环境与基线 & 现有依赖、最小输入和保存路线对应；今天的 CLI 重放与独立指标可查，不以旧 PASS 替代新验收。 \\
\addlinespace[4pt]
地图编辑 & 真实或合成场景可导入、修改、另存和重载；稳定 ID、坐标、洞、任务属性与原件证据不丢失。 \\
\addlinespace[4pt]
实际规划 & 编辑产生新修订；runner 产生新 job 和结果；scene\_\allowbreak{}revision、input\_\allowbreak{}hash、job\_\allowbreak{}id 与输出一致。 \\
\addlinespace[4pt]
核验和回放 & 独立重算客户恰一次、载荷不超 5、车辆不超 3 及合法线段；路线、成本和 checker 指向相同矩阵与代码。 \\
\addlinespace[4pt]
异常恢复 & 无解、搜索失败、超时、进程崩溃、取消、旧结果晚到与断网均保留准确状态和原始证据。 \\
\end{longtable}
\endgroup
% content block 4:23 h3
\Needspace{4\baselineskip}\subsubsection{性能目标尚待测量}
% content block 4:77 p
\bodytext{下列数值只是原型起始目标，尚无本项目性能跑分。测试记录设备、浏览器、分辨率、包版本、数据量和冷暖缓存；真实田块与固定 seed 合成压力集分别报告首次呈现、交互延迟、求解耗时、峰值内存和掉帧。}
% content block 4:12 table
\begingroup
\fontsize{9.6}{14.6}\selectfont
\setlength{\tabcolsep}{4pt}
\renewcommand{\arraystretch}{1.18}
\begin{longtable}{@{}>{\raggedright\arraybackslash}p{\dimexpr 0.201729\linewidth-4.8415pt\relax}>{\raggedright\arraybackslash}p{\dimexpr 0.446686\linewidth-10.7205pt\relax}>{\raggedright\arraybackslash}p{\dimexpr 0.351585\linewidth-8.4380pt\relax}@{}}
\toprule
{\sffamily\bfseries 场景} & {\sffamily\bfseries 拟定起始目标} & {\sffamily\bfseries 测试要求} \\
\midrule
\endfirsthead
\toprule
{\sffamily\bfseries 场景} & {\sffamily\bfseries 拟定起始目标} & {\sffamily\bfseries 测试要求} \\
\midrule
\endhead
\midrule
\multicolumn{3}{r}{\fontsize{8}{10}\selectfont\color{Muted}续下页}\\
\endfoot
\bottomrule
\endlastfoot
桌面普通场景 & 1 万可见顶点，平移缩放 45 fps；编辑反馈 p95 ≤ 200 ms。 & 同一数据集连续操作，记录设备与峰值内存。 \\
\addlinespace[4pt]
手机普通场景 & 2 千可见顶点，30 fps；编辑反馈 p95 ≤ 300 ms。 & 选定 Android、iOS 真机，覆盖触控冲突。 \\
\addlinespace[4pt]
大数据和栅格 & 按 1 千、1 万、10 万顶点及实际 COG 递增，找出降级点。 & 简化、切片、按需加载后重测；权威几何不变。 \\
\end{longtable}
\endgroup
% content block 4:78 p
\bodytext{正确性优先于帧率。坐标以参考控制点校验，阈值由精度和安全预算确定；矩阵、路径与服务事件独立重算。失败分类明确区分“证明无解”“搜索未找到”“超时”“执行异常”，不能把 INIT\_\allowbreak{}FAIL 或 repair\_\allowbreak{}failed 直接称为数学无解。}
% content block 4:24 h3
\Needspace{4\baselineskip}\subsubsection{扩展都需要通过自己的小试验}
% content block 4:79 p
\bodytext{整行与覆盖：先复用 P102、turn\_\allowbreak{}modes、Dubins 及 Fields2Cover 参考，做单任务类型的完整服务见证；以可信规则或传统规划器为简单基线。只有车辆与农具参数、任务数据及已测约束组合齐全，才能扩到多车、回补或动态执行；不支持的约束继续明确拒绝。}
% content block 4:80 p
\bodytext{部署与存储：先用响应式 Web、单 worker、SQLite 和有限区域包验证目标 Ubuntu 与手机，再考虑封装或独立离线。持续并发写或跨项目空间查询成为瓶颈时才评估 PostGIS；扩大只读影像不构成单独理由。多机器任务运营需要队列时考虑 Celery，复杂 CPU/\allowbreak{}GPU 任务图考虑 Ray，大表或数组数据流考虑 Dask；这些组件不替代幂等、重试和隔离。}
% content block 4:81 p
\bodytext{首轮交付包包含可编辑场景、冻结输入、原始结果、独立验算、依赖锁、结果 manifest 和回放入口。录屏证明操作过程，数值与完整见证证明算法和约束。同一待比较合同先冻结并通过自身运行资格；已有合格 CLI 可先开展机制实验，产品闭环并行验收，避免边改产品边改评价合同。}
% content block 4:13 source
\sourcetext{工程来源 T14、T19、T20、T21、T28，见技术来源附录。}
% content block 5:5 h1
\clearpage
\hypertarget{part_4}{}
\section{四 LLM理解 推荐与自动设计}
% content block 5:20 h2
\Needspace{6\baselineskip}\subsection{1  LLM理解 计划与受限执行}
% content block 5:25 h3
\Needspace{4\baselineskip}\subsubsection{在线任务编译与受限计划}
% content block 5:82 p
\bodytext{在线链为自然语言与场景→必要澄清→带来源的任务合同→语义检查与能力门控→性能证据→LLM候选调用计划→注册工具实际执行→独立验算。LLM须同时在场景理解与约束抽取、候选算法或组合计划推荐中留下可审计的实际作用，不能只复述选择器结论。[R15]}
% content block 5:83 p
\bodytext{计划只包含注册ID、参数预设、预算与回退。服务端拒绝未知工具、超额资源和任意表达式，不执行临时Python或shell。LLM原始提议、门控拒绝、实际执行ID、替换原因及求解回执全部保存；若最终采用确定性选择，明确记作混合策略，不能冒称LLM排序胜出。}
% content block 5:26 h3
\Needspace{4\baselineskip}\subsubsection{原意对齐先于路线合法}
% content block 5:84 p
\bodytext{checker能检查已给定合同，却抓不住合同漏掉的原意。例如“不得倒车”未被抽取，合法路线仍可能违反请求。因此保存原句→规范约束→检查条目；独立金标准核对否定、单位、数字、对象、硬软约束和目标。缺信息或矛盾时请求澄清，未知需求不默认为零。}
% content block 5:27 h3
\Needspace{4\baselineskip}\subsubsection{离线生成与演化}
% content block 5:85 p
\bodytext{冻结任务、函数接口、父代和训练侧反馈，生成受限评分/\allowbreak{}排序程序；可信内核掌握合法动作、真实成本和checker。原始回答、提取与执行源码、提示/\allowbreak{}父代哈希、模型、费用和失败均留存。静态检查、独立隔离、反例与开发评估通过后，才注册为在线可调用版本。}
% content block 5:86 p
\bodytext{P3保留思想与代码双表示和装箱比较；思想文字是算法说明，不能当模型内部推理。P5先复用BPP与CVRP/\allowbreak{}VRPTW，在同入口做生成、评估、保留更新、再生成、冻结测试和导出。评价后的决定须进入下一轮；若幸存父代仍是手写种子，应如实保留，不能强推弱子代制造演化成功。}
% content block 5:87 p
\bodytext{在线与离线共用任务、算法库、runner和结果包，权限和成本分开。较弱候选可留作研究，不强制部署；未有任何候选通过语义验算时，生成接入仍未完成。}
% content block 5:21 h2
\Needspace{6\baselineskip}\subsection{2  在线LLM对照与验收}
% content block 5:88 p
\bodytext{冻结同一库、输入权限和端到端预算，分开评价语义解析、工具计划和算法排序。借鉴OptiGuide的真实执行反馈、LLM-Modulo的外部验证、PICARD的结构化接口，但不直接引入其任意代码执行或声称语法正确就理解正确。[R15]}
% content block 5:13 table
\begingroup
\fontsize{9.6}{14.6}\selectfont
\setlength{\tabcolsep}{4pt}
\renewcommand{\arraystretch}{1.18}
\begin{longtable}{@{}>{\raggedright\arraybackslash}p{\dimexpr 0.158501\linewidth-3.8040pt\relax}>{\raggedright\arraybackslash}p{\dimexpr 0.458213\linewidth-10.9971pt\relax}>{\raggedright\arraybackslash}p{\dimexpr 0.383285\linewidth-9.1988pt\relax}@{}}
\toprule
{\sffamily\bfseries 策略} & {\sffamily\bfseries 保留或增加的机制} & {\sffamily\bfseries 回答的问题} \\
\midrule
\endfirsthead
\toprule
{\sffamily\bfseries 策略} & {\sffamily\bfseries 保留或增加的机制} & {\sffamily\bfseries 回答的问题} \\
\midrule
\endhead
\midrule
\multicolumn{3}{r}{\fontsize{8}{10}\selectfont\color{Muted}续下页}\\
\endfoot
\bottomrule
\endlastfoot
R0 无LLM & 结构化输入＋门控＋规则或统计选择。 & 共同正确合同上的推荐基线。 \\
\addlinespace[4pt]
R1 LLM基础 & LLM解释及推荐；保留全部执行安全限制，不合格提议记失败。 & 原始语义/\allowbreak{}推荐错误来自哪里。 \\
\addlinespace[4pt]
R2 受约束LLM & 增加schema、来源/\allowbreak{}语义检查、必要澄清和能力门控。 & 可靠任务编译带来什么增量。 \\
\addlinespace[4pt]
R3 性能辅助LLM & R2加训练侧性能、成本与范围证据，LLM提出受限选择或计划。 & 实测证据能否改善真实计划。 \\
\end{longtable}
\endgroup
% content block 5:89 p
\bodytext{另以同一金标准合同比较LLM合格集合内排序与统计规则，保持候选和预算相同；记录提议被确定性回退替换的比例。人工金标准输入不冒充免费语言解析，人工时间与信息分开报。 R2相对R1测的是组合效应；主张某条防线的独立作用时再拆消融。}
% content block 5:28 h3
\Needspace{4\baselineskip}\subsubsection{把解析与规划的损失拆开}
% content block 5:90 p
\bodytext{交叉比较人工金标准/\allowbreak{}LLM解析合同与固定SBS/\allowbreak{}实际选择四格；所有最终路线均按独立原意合同评分。它能辨认推荐收益、语言前端损失，以及解析错误是否特别伤害选择。终测按母场景和表达模板族留出，同义改写与最小对不跨分区。}
% content block 5:91 p
\bodytext{报告硬约束遗漏与虚增、整例语义正确、必要/\allowbreak{}多余澄清、原始计划合法、错误合同接受、按时合法完成、规划质量和全部延迟。全拒绝不能冒充可靠，格式合法不能冒充任务正确。去防线消融优先做离线影子回放；隐藏金标准checker仍保留，绝不为消融放开执行权限。}
% content block 5:92 p
\bodytext{首次语言交互、后续缓存快路径与人工等待分列；LLM排队、token、修复、门控、选择、求解、回退和checker全计。只有结果解释可读性提高时，就只报告该贡献。}
% content block 5:22 h2
\Needspace{6\baselineskip}\subsection{3  自动设计和选择的研究定位}
% content block 5:93 p
\bodytext{本项目最接近的已有研究不是单一LLM规划论文，而是自动配置、算法选择、程序搜索和组合构建的交叉。阅读的用途是缩小可成立的主张，并决定哪些模块直接借用、哪些必须适配。 本文SBS指训练侧选定、部署时固定使用的单一最佳方法；VBS指事后逐例挑最优的虚拟上界。selector是实际选择器，checker是独立验算器。}
% content block 5:29 h3
\Needspace{4\baselineskip}\subsubsection{采纳的共同骨架}
% content block 5:94 p
\bodytext{FunSearch、EoH与ReEvo支持“冻结模型生成程序、真实执行评分、用反馈继续搜索”的骨架。EoH的思想与代码双表示可对接老师主线；反思式反馈可作为受控方法对照。先固定一个生成实现和版本，再通过窄接口接URP，不同时维护多套重复框架。[R5、R6]}
% content block 5:95 p
\bodytext{ASlib和AutoFolio的价值在性能矩阵、特征费用、状态、分组及选择基线。首版实例推荐宜从廉价树、回归、kNN或规则开始。irace等自动配置挑战LLM是否只是有效地调参数；冻结程序检索挑战生成是否只是重用已有强父代。[R7、R8]}
% content block 5:30 h3
\Needspace{4\baselineskip}\subsubsection{直接限制原创性表述的工作}
% content block 5:96 p
\bodytext{Hydra已经将互补配置构建与实际算法选择相连；SATzilla成员贡献研究以删除成员后重建选择器来评估实际贡献。EoH-S关注互补启发式集合；LACE覆盖合同约束下的LLM组合构建，并提供rank-MILP选库。因而不能把“LLM＋互补库＋推荐”作为充分创新。[R9、R10、R11] DASH已有画像驱动的专长推荐；P142把选择器与执行费用纳入工作流优化，但其决策位置不同于执行前选求解器，需精确比较。[R21]}
% content block 5:97 p
\bodytext{建议检验的窄差异是：在有特征、探测、选择和回退费用的农业部署合同中，以分组交叉拟合的实际收益给候选信用，并识别它如何改变整个发现过程。实际成员边际、交叉拟合和成本计入各自都有先例；最终贡献必须来自具体机制、受控结果或有解释力的边界。}
% content block 5:31 h3
\Needspace{4\baselineskip}\subsubsection{不直接沿用的做法}
% content block 5:7 bullets
\begin{itemize}
\item 用逐实例事后最优的组合分数冒充可部署推荐。
\item 把源码哈希不同当作行为独立，把同一父代快照上的30个候选当30次独立发现。
\item 直接移植旧依赖和默认计时；原框架超时、失败、并列排名与许可证均须核对。
\item 只比较弱基线，或给LLM免费使用强父代与更多反馈。
\item 本地回查发现CALM附录一项基线行对应关系疑点，该项优势证据暂不采纳；不外推为整篇论文不可信。[R21]
\end{itemize}
% content block 5:23 h2
\Needspace{6\baselineskip}\subsection{4  扩库与推荐的四格主实验}
% content block 5:98 p
\bodytext{扩库子问题依次检验H1强库专长可被识别且净收益为正，H2生成优于同资产配置或检索，H3实际信用改善同池选库，H4该反馈改善完整发现。四项均是待检验假设；最后还需将冻结扩展库放回真实LLM在线链验证。 H1阴性不取消H2的功能验收；生成可能创造原池没有的能力，但须以有限新假设单独验证。}
% content block 5:99 p
\bodytext{在同一农业合同和总部署预算下，L0为强传统库，L1为训练侧冻结的L0加LLM候选。F为训练侧固定策略，S为实际选择器；损失越低越好。}
% content block 5:14 table
\begingroup
\fontsize{9.6}{14.6}\selectfont
\setlength{\tabcolsep}{4pt}
\renewcommand{\arraystretch}{1.18}
\begin{longtable}{@{}>{\raggedright\arraybackslash}p{\dimexpr 0.386167\linewidth-9.2680pt\relax}>{\raggedright\arraybackslash}p{\dimexpr 0.306916\linewidth-7.3660pt\relax}>{\raggedright\arraybackslash}p{\dimexpr 0.306916\linewidth-7.3660pt\relax}@{}}
\toprule
{\sffamily\bfseries 候选库} & {\sffamily\bfseries 固定策略 F} & {\sffamily\bfseries 实际选择 S} \\
\midrule
\endfirsthead
\toprule
{\sffamily\bfseries 候选库} & {\sffamily\bfseries 固定策略 F} & {\sffamily\bfseries 实际选择 S} \\
\midrule
\endhead
\midrule
\multicolumn{3}{r}{\fontsize{8}{10}\selectfont\color{Muted}续下页}\\
\endfoot
\bottomrule
\endlastfoot
L0 强传统库 & J00 & J01 \\
\addlinespace[4pt]
L1 加LLM候选 & J10 & J11 \\
\end{longtable}
\endgroup
% content block 5:100 p
\bodytext{报告旧库推荐J00−J01、扩库固定J00−J10、新库推荐J10−J11、同类选择政策下扩库J01−J11及整体J00−J11。讨论互补时再报(J10−J11)−(J00−J01)；交互为零不否定更强单算法的价值。选择器族、调参机会、信息和费用相同，两库各自训练。SBS还要面对同预算固定组合，合法率单报。}
% content block 5:32 h3
\Needspace{4\baselineskip}\subsubsection{挑战与同池比较}
% content block 5:101 p
\bodytext{Lcfg以同组件/\allowbreak{}参数空间做随机及自动配置，Lret以冻结强程序或父代检索；至少加入两者加S。强父代和历史反馈不能只免费给LLM。固定同一候选池、同库规模后，比较均值Top-k、EoH-S式CPI、LACE式rank-MILP、Hydra式增库与既有实际成员边际。[R9至R11]}
% content block 5:33 h3
\Needspace{4\baselineskip}\subsubsection{信用与发现分开验证}
% content block 5:102 p
\bodytext{内层按母组交叉拟合，估计L与L加候选在相同选择器和预算下的实际损失差，同时记录误选、合法率、尾部与费用。VBS改善而S不改善是未兑现机会；新成员普遍胜出则是单算法贡献。}
% content block 5:103 p
\bodytext{闭环只改变信用/\allowbreak{}反馈，其余模型、父代、接口与主预算固定；每次独立重启重新发现、冻结并到未见母组部署。同池有效不能证明生成分布改善。先做只换信用、只换选择器、移除成本及同预算多跑择优等最少消融，不把多个变化都归给一个机制。}
% content block 6:6 h1
\clearpage
\hypertarget{part_5}{}
\section{五 探索方向与最小实验}
% content block 6:24 h2
\Needspace{6\baselineskip}\subsection{1  探索组合与投入原则}
% content block 6:104 p
\bodytext{本章保留学习型路线改进、动态预算、LLM程序演化、开放模型后训练、生成多样性、agent执行与世界模型等研究空间。它们共同服务于农业任务的真实理解、合法执行、强算法库与实际推荐，但各自需要独立研究问题、数据和对照。}
% content block 6:105 p
\bodytext{投入顺序以主链为先：受约束LLM任务编译与注册工具实际执行，以及强传统库/\allowbreak{}LLM扩充库、固定策略/\allowbreak{}真实selector的四格归因。跨合同适应和运行中预算决策可优先作近端分支；其余方向保留资格与最小实验入口。这里的集中投入、暂缓或收缩均是研究安排建议，不表示用户取消某方向或正式课题能力。}
% content block 6:25 h2
\Needspace{6\baselineskip}\subsection{2  学习型路线改进与合同感知表示}
% content block 6:34 h3
\Needspace{4\baselineskip}\subsubsection{研究问题}
% content block 6:106 p
\bodytext{学习型路线方法分为三个可单独识别的角色：构造策略根据部分解生成下一服务动作；改进搜索从完整可行解中选择破坏、再插入或局部移动；决策小模型只判断哪个候选动作、子问题或昂贵评价值得执行。三者均作为算法库可替代成员，研究在真实通行、资源及方向信息完整后，是否仍能以更少总计算获得更好的合法路线。}
% content block 6:107 p
\bodytext{RRNCO类矩阵感知模型用于检验输入表示与路由决策；NDS作为学习破坏/\allowbreak{}再插入方向的前次调研候选入口；树、bandit、编辑价值头或图条件算子策略承担窄决策。应把补齐信息、网络表示、搜索动作和调度收益拆开，避免改换整个系统后无法归因。}
% content block 6:35 h3
\Needspace{4\baselineskip}\subsubsection{已有资产与事实边界}
% content block 6:108 p
\bodytext{已有NL124结果、精确参照与turn\_\allowbreak{}modes资格平台可直接复用，具体规模与支持域见“已有资产”及“转向与跨合同适应”。当前小合同的强P已达自身精确终值，学习实验应先确认时间到目标、合理规模或真实任务变化仍有改进空间。}
% content block 6:109 p
\bodytext{前次调研已留下NDS论文与代码入口，讨论用途为学习破坏/\allowbreak{}再插入，尚未复现。应明确标为“前次调研候选入口，版本及URP合同适配待核”，不能声称已在项目执行或已支持G/\allowbreak{}exact-K。RRNCO和URS已被原评议列为矩阵感知近邻，“加入真实G”本身不构成首次非欧氏建模。}
% content block 6:36 h3
\Needspace{4\baselineskip}\subsubsection{下一项最小实验}
% content block 6:110 p
\bodytext{先固定一个已验合同、传统初始化器、独立checker和总预算，绘制强传统方法的质量—成本曲线，确认仍有实用改进空间。构造、完整解改进与小模型决策分别做资格，不把任何一个局部成功解释为端到端学习系统已胜出。}
% content block 6:111 p
\bodytext{构造支线从同一合法状态编码开始，比较简单特征、矩阵编码及必要资源/\allowbreak{}模式状态；资格小例包括相同坐标而G不同、少车/\allowbreak{}空车、重复服务和全mask终态。三个不可拆需求各3、两辆容量5的车即使总需求9≤10仍不能完成，说明局部合法mask不保证最终可完成。损失比较另行固定输入：先同骨干BC、回报条件BC与离散IQL，再判断是否有必要增加历史或DT。}
% content block 6:112 p
\bodytext{NDS改进搜索支线先核定固定提交、模型接口、目标和动作含义，再以共同完整可行初解为起点，保持破坏规模、修复器、接受准则、停止条件和总预算一致，只替换破坏/\allowbreak{}再插入选择中的一个学习组件。先固定修复器比较学习破坏、随机破坏、距离或代价贡献规则、原生自适应选择；如果实现无法拆分，先做忠实整法复现，再明确整法结果不能单独归因于破坏策略。}
% content block 6:113 p
\bodytext{随后检验破坏与修复的交互：固定学习组件，仅替换共同的贪心/\allowbreak{}后悔插入或强传统修复器，记录可行恢复率、每次破坏—修复成本、期限内最优合法质量及重复启动方差。固定候选数和端到端墙钟分别报告。原生欧氏权重与目标合同若不一致，先适配/\allowbreak{}重训并计成本，不把缺G的旧权重作为唯一神经对手。}
% content block 6:114 p
\bodytext{决策小模型支线只预测少数昂贵修复、模式连接或子区域求解的收益；便宜且可精确计算的整数G即时delta继续直接算。固定执行器与候选动作，比较规则、树、bandit和小价值头，评价排序及错误选择的净损失；有稳定跨步效应后才进入动态预算分支。}
% content block 6:37 h3
\Needspace{4\baselineskip}\subsubsection{强简单对照}
% content block 6:115 p
\bodytext{所有角色共享同合同强PyVRP/\allowbreak{}HGS类求解器、正确mode-aware传统搜索和精确小例，不移除已胜出的强法。构造支线保留贪心构造、BC与回报条件BC；改进搜索支线保留固定邻域随机/\allowbreak{}贪心排序、随机与规则破坏、原生ALNS自适应权重、共同修复器下的多启动；决策小模型保留规则、树和简单bandit。}
% content block 6:116 p
\bodytext{若候选只调整温度、权重、破坏比例或move优先级，增加同空间随机搜索及irace。控制初始化、repair、局搜、信息权限、参数调优与端到端预算；允许传统基线使用与学习法相同的真实G和资源状态。}
% content block 6:38 h3
\Needspace{4\baselineskip}\subsubsection{数据条件}
% content block 6:117 p
\bodytext{构造数据至少记录完整状态、真实G、资源/\allowbreak{}模式、合法动作集、动作、后继、奖励、终止与费用。终局路线可回放为构造样本，但不能恢复原求解器的destroy动作、原来探索过的分支或其长程价值。}
% content block 6:118 p
\bodytext{改进搜索需记录破坏前完整解、被移除对象、破坏强度、再插入顺序或修复调用、接受/\allowbreak{}拒绝、后继解及全部成本；仅回放最终路线无法补出这些动作。NDS数据和权重资格须先核对原任务、G输入、车数、动作支持域与训练费用，暂不把前次调研入口算作已有轨迹。}
% content block 6:119 p
\bodytext{CDQAC的离散IQL可借损失与实现参照，其调度动作是工序×机器，buffer通过索引与并行环境数恢复后继，依赖同步等长采集。变长或异步VRP需显式保存episode、step、next\_\allowbreak{}state、合法集及终止原因，测试非法数据动作、超时截断和终态。其低质数据结果受密集目标一致奖励、完整等长轨迹等条件约束，不能外推为随机数据普遍更好。}
% content block 6:120 p
\bodytext{新母田块/\allowbreak{}母图按组留出，旧NL124和反复查看的turn根图继续作开发与资格集。训练、建图、GPU传输、推断、局搜及checker均计费；公开权重需重新核对欧氏目标、车数与规模支持域。}
% content block 6:39 h3
\Needspace{4\baselineskip}\subsubsection{近邻与研究增量}
% content block 6:121 p
\bodytext{RRNCO/\allowbreak{}URS限制非欧氏表示的新颖性主张；NDS使学习破坏/\allowbreak{}再插入成为必须面对的近邻，现阶段仍是版本与URP合同适配待核的候选入口。DR-ALNS已有算子、破坏强度和接受温度学习，不能将这些角色第一次接在一起宣称为新机制。}
% content block 6:122 p
\bodytext{CDQAC提供离散离线学习参照，RL4CO提供环境接口。项目增量应落在农业合同需要的状态、可完成性、同成本动作决策、修复交互或跨母场景表现；使用GNN、IQL、NDS或历史窗口本身不是贡献。}
% content block 6:40 h3
\Needspace{4\baselineskip}\subsubsection{保留与收缩条件}
% content block 6:123 p
\bodytext{若在未见母场景、同信息同成本下仍提高合法率或anytime质量，且简单规则无法解释增益，保留该学习成员并进入统一算法库四格实验；模型可能仅在一个合同或预算区间有用，可按支持域注册。}
% content block 6:124 p
\bodytext{若正确传统状态扩展已解释全部收益，或新增模型只击败缺信息/\allowbreak{}违规版本，应将结论收缩为可靠建模。若强求解器已精确解决目标、推断开销吞掉改进或泛化失败，不扩大网络规模；这不取消学习型路线方向，也不删除正式推荐功能。}
% content block 6:14 source
\sourcetext{重要来源：[ER51]、[ER52]、[ER01]、[ER02]、[ER03]、[ER04]、[ER05]、[ER47]、[ER50]。项目历史事实另见[R4]；ER编号对应文末探索方法来源。}
% content block 6:26 h2
\Needspace{6\baselineskip}\subsection{3  动态调度与运行中预算决策}
% content block 6:41 h3
\Needspace{4\baselineskip}\subsubsection{研究问题}
% content block 6:125 p
\bodytext{动态分支研究下一次计算是否值得：在已观察的运行进展与剩余预算下，继续当前程序、重启、切换求解器、调用哪个算子或提高评价深度，哪项决策能改善期限内的合法结果？实例前选solver、运行中选动作、发现时选父代/\allowbreak{}评价、固定池选库是四个不同接口，不能用其中一层的成功替另一层补证。}
% content block 6:42 h3
\Needspace{4\baselineskip}\subsubsection{已有资产与事实边界}
% content block 6:126 p
\bodytext{已有BPP、VRPTW、CVRP性能矩阵可以诊断固定库与静态选择；仅有终局分数不足以学习序贯控制。原评议记录PyVRP观察器缺动作和逐步路线，不能从统计进展补造完整状态—动作轨迹。当前runner是否支持真实继续与切换仍需逐算法核验。}
% content block 6:127 p
\bodytext{NAS另有8架构×5发现seed和5确认seed，共80条完整训练轨迹、1600行epoch记录；2 epoch为同次20 epoch训练的前缀，两种保真度选中同架构。它可作等质省计算的记录与多保真诊断资产，不能当农业控制监督；checkpoint真实恢复语义仍需核查。}
% content block 6:43 h3
\Needspace{4\baselineskip}\subsubsection{下一项最小实验}
% content block 6:128 p
\bodytext{先选择一种合同和少量强程序，记录完整anytime曲线与合法incumbent，辨别是否出现可重复且部署时可预测的预算—质量差异。区分三种变化轴：一个程序随运行时间改善、生成谱系随代数改善、同一组件随搭档变化才有用。它们各需固定另外两轴。}
% content block 6:129 p
\bodytext{只有存在实际空间，再进行可复位的同状态分支：固定后续策略、剩余预算和环境条件，对继续、重启及少数昂贵动作做成对诊断。随后只改变低频控制策略，固定算法库、总预算与回退。若即时delta可直接用整数G计算，优先预测深修复、高成本几何或长程组合的价值，避免用网络替代便宜加减法。}
% content block 6:44 h3
\Needspace{4\baselineskip}\subsubsection{强简单对照}
% content block 6:130 p
\bodytext{强单算法全预算、训练侧静态selector、固定时间切分与顺序、预定重启、原生ALNS自适应权重、停滞规则、树与简单bandit。先挑战这些策略，再增加长程RL或价值备份。}
% content block 6:45 h3
\Needspace{4\baselineskip}\subsubsection{数据条件}
% content block 6:131 p
\bodytext{状态需包含剩余预算、合法incumbent、已观察进展、必要的搜索/\allowbreak{}资源状态；记录动作、动作可用性、后继、终止、超时与费用。日志中低平均收益可能来自算子总在困难状态被调用，不能直接当作因果效应。仅用离线日志时必须说明状态充分性、行为覆盖与分布偏移。}
% content block 6:132 p
\bodytext{probe、状态复制、保存/\allowbreak{}加载、切换、冷启动、推断和checker全计费。无法resume的程序按冷启动处理，不能把A从零十秒后的好成绩与先跑B一秒的曲线拼接成免费包络。并行墙钟与CPU/\allowbreak{}GPU总时分别记录。}
% content block 6:46 h3
\Needspace{4\baselineskip}\subsubsection{近邻与研究增量}
% content block 6:133 p
\bodytext{DR-ALNS已学习算子、破坏强度和接受温度，Learning to Delegate已选择交给传统求解器的子问题。MuEvo主要处理协作组件及可逆生命周期，其FE不含所有探测、生成和修复成本，并不自动解决任意程序慢启动删失。新意需要来自明确的状态信息、可恢复决策或净成本机制。}
% content block 6:47 h3
\Needspace{4\baselineskip}\subsubsection{保留与收缩条件}
% content block 6:134 p
\bodytext{中间状态相对静态特征带来可重复预测增量，且决策收益扣除控制与恢复成本后仍超出预定实用阈值，才增加策略复杂度。导航碰撞检查、覆盖连接、VRP子区域修复可复用这一评价原则，但动作网络和标签不自动通用。}
% content block 6:135 p
\bodytext{没有排名反转、只有人工延迟才能产生信号、固定调度实用等效或高频回调吞掉节约，就保留静态/\allowbreak{}低频规则。未收集够动作覆盖时先补诊断，不能把缺数据等同于RL方向无效，也不能事后改变预算定义制造收益。}
% content block 6:15 source
\sourcetext{重要来源：[ER05]、[ER06]、[ER07]、[ER08]、[ER48]。项目历史事实另见[R4]；ER编号对应文末探索方法来源。}
% content block 6:27 h2
\Needspace{6\baselineskip}\subsection{4  LLM生成 修复与可迁移程序演化}
% content block 6:48 h3
\Needspace{4\baselineskip}\subsubsection{研究问题}
% content block 6:136 p
\bodytext{保持冻结LLM也能研究自动算法设计。核心问题是同一父代和修改预算下，标量分数、结构化失败定位、可执行反例或历史检索，哪些信息能提高合法改进的获得率，并将收益保留到未见母实例、错误族和代码家系？合同适配与启发式优化应分别评价。}
% content block 6:49 h3
\Needspace{4\baselineskip}\subsubsection{已有资产与事实边界}
% content block 6:137 p
\bodytext{P2批次60槽生成档案已有提示、父代、回答、CLI回执及哈希，分阶段43完整评分、7静态拒绝、10生成失败；EOH共享冻结父代快照。P1、P2是两批生成记录，并非连续两代或独立发现重复。生成runner负责请求留痕，不负责评分、父代更新或入库。}
% content block 6:138 p
\bodytext{另有紧凑DAG的22次物理请求和真实选择记录；10请求含FF/\allowbreak{}BF父程序，14次父代引用均来自FF/\allowbreak{}BF，没有生成后代被选作后续父代，最终均选FF。15候选结构拒绝、7进入四例评价、仅4个完整有效fitness。它是有价值的局部编辑诊断，不能与完整Python程序档案混成一条扩张谱系。}
% content block 6:139 p
\bodytext{新20序列EOH负结果及IID缺失比较沿用“已知失效”口径，不因后续设计而覆盖原始结果。}
% content block 6:50 h3
\Needspace{4\baselineskip}\subsubsection{下一项最小实验}
% content block 6:140 p
\bodytext{固定模型、父代、可编辑函数、候选数和旧能力保护集，仅替换反馈。比较标量分数、结构化定位、完整失败实例、保语义缩减反例，并加入不需要修改的任务及同长度无关反馈。分别用固定token和固定评价核秒报告，不把提取代码更方便或更多反馈token当规划知识。}
% content block 6:141 p
\bodytext{优先从turn\_\allowbreak{}modes的窄接口试验：合同适配只修车辆数、depot或成本接口；启发式迁移则保持服务与车队语义，改变转向成本/\allowbreak{}禁转，仅提议顺序或邻域优先级，固定模式DP与checker。比较隐藏合同通过、旧能力回归、合法质量及适应全成本。}
% content block 6:142 p
\bodytext{迁移反例必须真的区分合同。NL124有12个单位需求、Q=5、K=3，在单趟无补给条件下，至多3车与恰3车可行域相同，不能用于证明该迁移。可先使用2个单位需求、Q=2、K=2的见证；缩减反例须保留输入合法、相同失败谓词及目标可行性。}
% content block 6:51 h3
\Needspace{4\baselineskip}\subsubsection{强简单对照}
% content block 6:143 p
\bodytext{源程序零样本、不改代码而调用正确新evaluate的强mode-aware搜索、目标从零生成、独立重采样、冻结检索、手写一次适配、同空间随机搜索/\allowbreak{}irace，以及成熟修复循环。表达接口可比较自由Python、窄API和typed DSL，但空间内公平对照与对外部强手写法的能力差距分开报告。}
% content block 6:52 h3
\Needspace{4\baselineskip}\subsubsection{数据条件}
% content block 6:144 p
\bodytext{完整编辑单元包含经确认的旧/\allowbreak{}新合同、父代及家系、允许补丁、原始反馈与失败见证、所有候选、保护集、分阶段结果和全成本。程序哈希与行为/\allowbreak{}家系去重分开；测试反馈不得进入检索档案。词法、AST与行为指纹检索共用能力过滤，行为测试另收费。}
% content block 6:145 p
\bodytext{小型合同试验先识别机制，完整发现结论再做独立流程重启、开发侧选择和未触碰终测。原turn的12根图和反复查看的数据作为开发集，新增母图独立留出。父代选择回放只改变旧候选选择分布，不是新后代实验。}
% content block 6:53 h3
\Needspace{4\baselineskip}\subsubsection{近邻与研究增量}
% content block 6:146 p
\bodytext{EoH已有思想和代码共同进化，AlgoEvo已有诊断经验与技能检索，AHD Agent已有工具诊断及训练后生成策略，“生成—评价—反思”本身不是新机制。ChopChop提供解码时语义约束入口，但普通闭源chat接口未必有逐token控制。LLM-CEGIS-Repair虽有失败反馈和带洞草图，其已核代码没有硬验证洞外不变，并存在依据stdout而非严格退出成功门控的风险；这是静态审计问题，不是已触发复现。应独立采用基础设施失败标unknown、失败即拒绝通过的checker。}
% content block 6:54 h3
\Needspace{4\baselineskip}\subsubsection{保留与收缩条件}
% content block 6:147 p
\bodytext{反例反馈若在未见错误族上降低修复成本且不增加回归，可形成程序可靠性贡献；若合法程序进一步在独立母场景上胜过强搜索/\allowbreak{}配置，才增加算法优化或迁移知识主张。生成候选必须经离线准入再进入在线注册库。}
% content block 6:148 p
\bodytext{若只换正确cost即可适配、手写适配/\allowbreak{}简单检索同效、收益由免费父代或额外搜索解释，就收缩相应机制主张。成功调用现成DP属于可靠适配，不称新规划算法；固定DP和邻域、只学优先级也不称新邻域发现。}
% content block 6:16 source
\sourcetext{重要来源：[ER09]、[ER10]、[ER11]、[ER12]、[ER13]、[ER14]、[ER15]、[ER16]、[ER50]。项目历史事实另见[R4]；ER编号对应文末探索方法来源。}
% content block 6:28 h2
\Needspace{6\baselineskip}\subsection{5  开放LLM后训练与小顾问模型}
% content block 6:55 h3
\Needspace{4\baselineskip}\subsubsection{研究问题}
% content block 6:149 p
\bodytext{后训练的研究问题是同一批任务、程序和反馈，写入参数是否比冻结检索或额外搜索更有效，而不是接口能否微调。完整程序生成、局部编辑和短策略顾问是不同训练对象，应先选一个；普通GLM API改提示或继续调用不属于权重后训练。}
% content block 6:56 h3
\Needspace{4\baselineskip}\subsubsection{已有资产与事实边界}
% content block 6:150 p
\bodytext{当前43个评分程序共用同序五条设计输入。IID的19个不同程序共享1个prompt，形成156个严格偏好对；EOH的24程序属于20个prompt，其中6程序、3提示形成3对。159条偏好边中98.1\%来自同一IID提示，不能视为159份独立任务，也不能凭代码哈希不同推断家系独立。跨臂提示/\allowbreak{}程序交集尚未核。}
% content block 6:151 p
\bodytext{这些记录可支持单BPP合同的标签与工程资格诊断，尚无可迁移参数信号或农业迁移效果。终局求解表本身不含真实编辑链；无代码的请求失败不能作代码负例，准入拒绝和服务异常分开。公开EvoTune、CALM、Guidance-TTT曾做固定源码核查，不等于项目已安装、训练或运行；本地单个calm.py也不是CALM训练系统。}
% content block 6:152 p
\bodytext{已有硬件类别记录不证明能训练所选模型。Qwen3.5-4B与Gemma 4 12B官方卡可作开放候选入口，许可、显存、兼容性、mask、吞吐和重载均需实测，不能据型号承诺工期。}
% content block 6:57 h3
\Needspace{4\baselineskip}\subsubsection{下一项最小实验}
% content block 6:153 p
\bodytext{先完成一个开放模型的冻结生成、数据mask、一次小规模适配、adapter保存/\allowbreak{}重载及独立评价资格。若借EvoTune，固定683b245251e9的src/\allowbreak{}experiments/\allowbreak{}main.py，并比较none.yaml与dpo.yaml；其每轮从基座重新训练LoRA，不应描述为接着上一adapter持续更新。}
% content block 6:154 p
\bodytext{最小有效比较为原基座冻结、原基座加同档案检索、一次LoRA/\allowbreak{}DPO加同检索，另给冻结组等资源追加搜索。共同未见父代或生成任务用于识别权重能力，开发侧锁定checkpoint与程序；若声称DPO负例/\allowbreak{}偏好机制有独立价值，再加入同数据高质程序SFT。一次适配后多次搜索只评价该checkpoint，训练方法稳定性需在更外层重复数据采样、适配与checkpoint选择。}
% content block 6:155 p
\bodytext{小顾问可作为更窄入口：固定强coder、同父代和实现预算，比较规则建议、冻结顾问、廉价排序器与可训练顾问，绑定建议、实际实现、失败和judge结果。若目标是对该coder的实际效用，末端成绩可作标签；只有主张独立于coder的算法思想时，才需重复实现或跨coder检查。}
% content block 6:58 h3
\Needspace{4\baselineskip}\subsubsection{强简单对照}
% content block 6:156 p
\bodytext{冻结强模型、同档案检索、等资源追加搜索、同数据SFT；小顾问使用规则、冻结顾问及小排序器。所有组共享开发资料权限和逐状态上下文，并报告实际曝光量，不能把训练组读过更多历史、继承更好种群的收益全归给权重。}
% content block 6:59 h3
\Needspace{4\baselineskip}\subsubsection{数据条件}
% content block 6:157 p
\bodytext{按合同、评价面板、prompt、程序和谱系整理档案，并以连通关系分组留出；对全pair使用prompt/\allowbreak{}程序采样或加权敏感性检查，ties、测量不确定和服务故障不强行排序。标签可在训练侧重评，终测保持未触碰。EvoTune/\allowbreak{}DAR-DPO限制一次采样中程序复用，全pair并非不合法，但不能声称已等价复现其协议。}
% content block 6:158 p
\bodytext{SFT需可信目标，偏好训练需共同任务下可比优劣，局部编辑另需共同编辑状态和可见信息；在线RL需要足够动作采样与执行预算。若同prompt的全部reward减去同一父代常数后再组内中心化，该常数完全抵消，应先核对advantage是否实际改变。}
% content block 6:159 p
\bodytext{完整成本包括标签、失败、训练、开发侧选checkpoint、推断与搜索。固定checkpoint评价与外层训练重复分开；若每次未来设计任务净节约ΔC>0，单一成本维度回本需任务数超过C\_\allowbreak{}train/\allowbreak{}ΔC，不能把一个算法服务很多规划请求当成很多次设计节约。}
% content block 6:60 h3
\Needspace{4\baselineskip}\subsubsection{近邻与研究增量}
% content block 6:160 p
\bodytext{EvoTune已有迭代偏好优化；CALM正式版已有DPO/\allowbreak{}GRPO与时间分解，不能沿用“未给训练耗时”的旧说法。其CVRP宿主从坐标重算欧氏距离、容量50、内部seed42、所核迭代常数100且无固定K，迁移到G/\allowbreak{}exact-K须重写合同入口、完整路线输出与独立验算，许可和依赖另核。}
% content block 6:161 p
\bodytext{Guidance-TTT和PACEvolve++已有小顾问加强执行器，角色拆分本身不是新增机制。Guidance-TTT所核默认1步、2×4 rollout不同于论文级30步、8×16，测试源码存在不等于项目通过。From Search to Signal已有结果信用、种群与权重拆分，其三seed主队列并未普遍胜出，预算锚也不等于逐次精确同墙钟。AHD-Agent所核仓树仅README/\allowbreak{}许可/\allowbreak{}图像，不能当现成训练代码。}
% content block 6:61 h3
\Needspace{4\baselineskip}\subsubsection{保留与收缩条件}
% content block 6:162 p
\bodytext{当新任务/\allowbreak{}家系上的合法改进率、编辑成本或设计总成本稳定超过冻结检索与追加搜索，再扩大数据和模型。训练尚未回本也可形成机制研究，但要明确效益对象和适用范围。}
% content block 6:163 p
\bodytext{极少独立提示组只做资格诊断，不立刻作训练方法泛化结论。若收益在同档案强基线后消失，或由种群/\allowbreak{}曝光/\allowbreak{}更多搜索解释，就暂停额外训练投入并保留冻结路线；这是资源安排与主张收缩，不是用户取消后训练方向。}
% content block 6:17 source
\sourcetext{重要来源：[ER17]、[ER18]、[ER19]、[ER20]、[ER21]、[ER22]、[ER23]、[ER24]、[ER25]、[ER26]、[ER27]、[ER28]、[ER29]。项目历史事实另见[R4]；ER编号对应文末探索方法来源。}
% content block 6:29 h2
\Needspace{6\baselineskip}\subsection{6  生成多样性与变化后的恢复价值}
% content block 6:62 h3
\Needspace{4\baselineskip}\subsubsection{研究问题}
% content block 6:164 p
\bodytext{GFlowNet、扩散或掩码重构的价值应体现在同生成和后处理预算下覆盖不同有用结构，提高达标概率或环境变化后的恢复能力。原始采样熵、路线表示个数、真正不同的方案以及后处理后的实际效用必须分开。}
% content block 6:63 h3
\Needspace{4\baselineskip}\subsubsection{已有资产与事实边界}
% content block 6:165 p
\bodytext{当前对称G、同质3车且路线方向无业务含义时，同一方案通常有3!×2³=48种串行表示；48条不同轨迹可能仍是一个方案。非对称G不能反向去重，异质车不能随意换车号。等价关系必须由任务合同决定，不能统一套用。}
% content block 6:166 p
\bodytext{GFACS公开固定源码已有GNN边特征、ACO计价和SWAP*距离矩阵链，不能写成只能用坐标；当前对称G可借入口。有向G需审计depot双向特征、对称信息素和反向等价；原生构造无一般exact-K，局搜上限允许空路线，必须统一采样、重评分、后处理与checker的支持集。}
% content block 6:64 h3
\Needspace{4\baselineskip}\subsubsection{下一项最小实验}
% content block 6:167 p
\bodytext{先用传统候选档案检验多样性是否有用途：变化发生前冻结单最佳解、质量匹配的传统多启动档案和学习候选库，规范化去重后控制候选数和原始质量，统一后处理。先测静态达标概率，再用预登记的有限边成本扰动G→G′测热启动质量、恢复时间和错选损失。}
% content block 6:168 p
\bodytext{边成本改变而容量/\allowbreak{}K不变时，旧解通常仍满足资源合同，可重评；容量、车数或通行结构变化另测可行恢复，不能混作同一终点。学习机制最小实验固定encoder、合法解码、reward、repair和局搜，仅比较TB与一种替代损失，或一种预定义结构覆盖机制。候选数配平回答样本效率，端到端墙钟配平回答部署效率。}
% content block 6:65 h3
\Needspace{4\baselineskip}\subsubsection{强简单对照}
% content block 6:169 p
\bodytext{传统多启动、无学习ACO、同骨干策略梯度、简单采样与多样性档案，以及单最佳解热启动。引入扩散时固定解码能力和局搜预算，避免把更强后处理归给去噪网络。}
% content block 6:66 h3
\Needspace{4\baselineskip}\subsubsection{数据条件}
% content block 6:170 p
\bodytext{保存原始路线、合同规范化表示、质量、结构距离、候选来源和后处理全部轨迹。若后处理随机，一次相同终点不足以证明所有样本属于同一完整吸引域。候选库必须在变化发生前构建，选择器只能看当时允许的信息。}
% content block 6:171 p
\bodytext{资格例包括重复客户、少车、空路线、错报成本。GFACS原validate\_\allowbreak{}route的集合覆盖不能排除重复服务，NLS每候选含三段局搜，公开测试计时遗漏GNN前向等步骤；本项目须把建图、推断、传输、全部局搜和checker纳入部署账。训练与标签费用另列，“无近优标签”不等于无计算。}
% content block 6:67 h3
\Needspace{4\baselineskip}\subsubsection{近邻与研究增量}
% content block 6:172 p
\bodytext{GFACS已把学习分布接入ACO且处理相关反向概率；f-Trajectory Balance提供损失替换入口，CAM降低近优标签依赖。新增研究应落在合同一致的支持集、质量匹配后的结构覆盖或扰动恢复，不能把GFlowNet接搜索、更多轨迹或对称去重单独宣称创新。}
% content block 6:68 h3
\Needspace{4\baselineskip}\subsubsection{保留与收缩条件}
% content block 6:173 p
\bodytext{等价去重、原始质量匹配和统一后处理后，学习候选库仍在新母任务或预定变化上改善达标/\allowbreak{}恢复与净成本，才支持多样性机制。导航通道、覆盖条带方向与区块连接、车辆客户分组可形成真实差异，但只迁移评价原则，不直接共用等价关系。}
% content block 6:174 p
\bodytext{若优势经规范化消失、统一局搜抹平差异、只是产生更多低质量方案，或传统多启动同墙钟更优，就保留廉价档案并收缩生成分布的部署主张。忠实复现及失效边界可保留为成果，不必取消方向。}
% content block 6:18 source
\sourcetext{重要来源：[ER30]、[ER31]、[ER32]、[ER33]、[ER34]、[ER35]。项目历史事实另见[R4]；ER编号对应文末探索方法来源。}
% content block 6:30 h2
\Needspace{6\baselineskip}\subsection{7  Agent与harness的受约束执行决策}
% content block 6:69 h3
\Needspace{4\baselineskip}\subsubsection{研究问题}
% content block 6:175 p
\bodytext{Harness是模型外的工具、沙箱、版本与执行循环，agent在其中决定下一步调用。方法问题应聚焦何时诊断、重新生成、换父代、增加测试、请求澄清或终止，以及这些决策是否在同预算下提高任务忠实完成率。增加工具和多个角色本身首先是工程实现。}
% content block 6:70 h3
\Needspace{4\baselineskip}\subsubsection{已有资产与事实边界}
% content block 6:176 p
\bodytext{已有真实生成、选择、CLI证据链、BPP面板与历史地图原型；当前产品接入缺口和任务状态要求见技术部分。固定CLI可先做机制实验，前后端全接通不是科研比较前置。}
% content block 6:177 p
\bodytext{OpenLoopX/\allowbreak{}LLM4AD Next可作为任务构建、evaluator、代码块建议及EoH、ReEvo、MCTS-AHD的后端来源，README中Impl与Skill不同、FunSearch仍标待实现；依赖、Python版本和合同需实测，仓库存在不等于能直接运行。无需因此替换既有面板，固定CLI也能先做机制试验。}
% content block 6:71 h3
\Needspace{4\baselineskip}\subsubsection{下一项最小实验}
% content block 6:178 p
\bodytext{冻结工具集合、版本、checker、自然故障集合与总预算，仅改变何时调用。对比固定状态机、强规则、廉价bandit和学习agent；先保留结构化任务输入，再单独增加自然语言入口，避免把语义解析与调度收益混合。}
% content block 6:179 p
\bodytext{故障试验可含超时、重复提交、陈旧地图、无效代码、错误单位和缺失约束；自然失败与人工注入分别报告。评价任务忠实完成、错误终止、无谓调用、恢复费用和最终规划质量，不只统计工具调用成功。恢复动作涉及地图或合同变化时，必须更新任务版本并使旧成本、旧结果及旧证据按规则失效。}
% content block 6:72 h3
\Needspace{4\baselineskip}\subsubsection{强简单对照}
% content block 6:180 p
\bodytext{固定状态机、强规则、标准重试/\allowbreak{}幂等处理、人工设计的最小澄清规则及简单bandit；相同模型的无学习固定流程。两角色或多agent不是独立性对照，需单独核原意、形式约束与几何真值。}
% content block 6:73 h3
\Needspace{4\baselineskip}\subsubsection{数据条件}
% content block 6:181 p
\bodytext{保存原用户意图、必要澄清与确认、任务合同及来源、工具计划、实际执行ID、工具回执、状态转移、终止理由、失败见证和总费用。离线代码修改与在线工具调用采用不同权限边界：普通在线计划只调用已准入注册版本，地图编辑不自动授权运行新代码。}
% content block 6:182 p
\bodytext{用户意图与可视编辑需要人工金标和真实修改任务；把“纠正理解”“改变任务”“更新偏好”“表达不确定”分开。公开XML/\allowbreak{}XSD不包含项目exact-K、容量、全行服务或机具扫掠，不能直接当中文农业合同金标。规则回放可验接口，节省用户时间和减少漏错需要真实参与者证据。}
% content block 6:74 h3
\Needspace{4\baselineskip}\subsubsection{近邻与研究增量}
% content block 6:183 p
\bodytext{AgriAgent已有按能力合同检索、组合工具及生成测试注册；As You Wish已有语言转XML、独立LTL生成和SPIN反馈，农业“语言+工具+验证”不是空白。其300英尺请求与约285×275英尺田块例说明形式验证仍需几何元约束，但不是实车越界事故，也不能只凭两条边长排除300英尺斜向路径。项目可研究原意、形式合同、实际几何的交叉验算。}
% content block 6:75 h3
\Needspace{4\baselineskip}\subsubsection{保留与收缩条件}
% content block 6:184 p
\bodytext{若在未见故障组合和真实任务上，学习调用策略比规则降低错误终止/\allowbreak{}恢复成本，且任务忠实与最终质量不退化，可研究策略迁移。即使性能推荐无显著净收益，也继续保留正式课题要求的LLM理解、推荐与实际调用，并如实报告其功能或可靠性价值。}
% content block 6:185 p
\bodytext{规则实用等效时保留规则，不追加通用多agent扩张、自改评估器或自动扩权。纯恢复成功支持系统可靠性，不自动证明优化算法更强；漂亮面板、工具数量和角色分工不作为方法增量。}
% content block 6:19 source
\sourcetext{重要来源：[ER36]、[ER37]、[ER38]、[ER39]。项目历史事实另见[R4]；ER编号对应文末探索方法来源。}
% content block 6:31 h2
\Needspace{6\baselineskip}\subsection{8  世界模型与昂贵后果评价的代理}
% content block 6:76 h3
\Needspace{4\baselineskip}\subsubsection{研究问题}
% content block 6:186 p
\bodytext{世界模型保留为条件性方向，回答给定状态和行动之后的未知后果，或以可验证代理替代昂贵高保真评价。当前静态G已知时，再学习同一距离不会自动增加决策信息；显式禁区可用多边形表达，只有滑移、能耗、载荷影响、土壤压实或高成本模拟确实改变计划选择时，才需要扩展模型。}
% content block 6:77 h3
\Needspace{4\baselineskip}\subsubsection{已有资产与事实边界}
% content block 6:187 p
\bodytext{现有坐标、图、容量与运动学组件支持显式几何和任务验算，不构成未知轮地动力学数据。RELLIS-3D full-stack包包含cmd\_\allowbreak{}vel、轮式里程计、惯导/\allowbreak{}GPS及电流电压；不能依据主要同步图像与点云的synced包说它只有感知标签。cmd\_\allowbreak{}vel是速度控制层，不是轮胎扭矩，也不能提供未采取动作的任意反事实。}
% content block 6:188 p
\bodytext{TartanDrive 2.0有公开七小时、最高15m/\allowbreak{}s的离路数据及下载工具，高速ATV不是低速带农具工况。Chrono地形模型可作模拟起点，默认土壤参数不是项目标定。原材料只确认公开入口，数据已取得、已同步、已校准和项目已训练必须另行登记。}
% content block 6:78 h3
\Needspace{4\baselineskip}\subsubsection{下一项最小实验}
% content block 6:189 p
\bodytext{固定同一批传统算法产生的合法候选计划，先在独立观测或经过校准的模拟器上比较静态成本、运动学加保守误差界、小残差模型与学习模型。首要终点是计划排序、错选损失、危险阈值附近的尾部误判、校准和全成本，之后才检验模型是否帮助LLM生成更稳健算法。}
% content block 6:190 p
\bodytext{未知后果路线先做行动对齐与必要的同状态不同动作检查；昂贵已知评估路线先验证代理误差以及独立高保真回查。超出行为覆盖时保留拒判或保守策略。模型预测画面逼真、视频相似或自身模拟得分更高均不能替代下游决策证据。}
% content block 6:79 h3
\Needspace{4\baselineskip}\subsubsection{强简单对照}
% content block 6:191 p
\bodytext{静态真实成本、显式运动学、保守上界、小线性/\allowbreak{}树残差模型及等成本追加精确评测。对比必须共用候选计划、可用观测、回查预算和风险约束，不能因世界模型多看信息而归功于模型结构。}
% content block 6:80 h3
\Needspace{4\baselineskip}\subsubsection{数据条件}
% content block 6:192 p
\bodytext{状态—控制—外部变化—后继必须按控制层级和时间戳对齐，并记录载荷、定位误差、传感质量及场景支持范围。行动后果训练需要足够行为覆盖；农业压实、作物损伤和车辙可能不可逆，不预设正反动作相互抵消。}
% content block 6:193 p
\bodytext{RELLIS官方60秒full-stack示例约4.2GB，取得和同步有真实成本；仓库标CC BY-NC-SA 3.0，Tartan代码MIT不自动覆盖全部外部数据。采集、标定、训练、推断、数据传输、模拟维护和独立高保真回查全部记账。没有工况校准时可报告参数敏感性或数字试验场，不能称数字孪生或真实部署安全。}
% content block 6:81 h3
\Needspace{4\baselineskip}\subsubsection{近邻与研究增量}
% content block 6:194 p
\bodytext{DINO-WM依赖动作对齐离线数据；JEPA-WMs指出多步预测实现修正会影响规划，即使画面相近；CoCo研究忽略行动影响的捷径。项目增量要由农业工况、行动可辨别性、风险校准或昂贵评估替代的净价值支持，不以“接入世界模型”单独作为贡献。}
% content block 6:82 h3
\Needspace{4\baselineskip}\subsubsection{保留与收缩条件}
% content block 6:195 p
\bodytext{只有模型在独立真实观测或已声明参数范围的高保真回查中，显著减少错选与尾部风险，且强简单模型/\allowbreak{}额外精确评测不能同成本达到，才纳入主线。预测改进与后续LLM发现收益分别识别。}
% content block 6:196 p
\bodytext{若简单残差或保守界已同效、未知因素不改变决策、行为覆盖不足或收益仅在模型自产环境成立，就暂缓复杂世界模型主投入，保留数据与模拟资格研究；这不是取消该探索方向。没有行动数据也不应靠生成视频补造物理真值。}
% content block 6:20 source
\sourcetext{重要来源：[ER40]、[ER41]、[ER42]、[ER43]、[ER44]、[ER45]、[ER46]。项目历史事实另见[R4]；ER编号对应文末探索方法来源。}
% content block 6:32 h2
\Needspace{6\baselineskip}\subsection{9  可视编辑与澄清的研究实验}
% content block 6:197 p
\bodytext{已有地图编辑原型、TaskSpec与独立checker可以组成交互实验；界面工程是否跑通和用户是否更少犯错分别验收。主问题是可视编辑、来源展示和有针对性的澄清，能否提高最终任务忠实度，而非仅提高满意度。}
% content block 6:83 h3
\Needspace{4\baselineskip}\subsubsection{下一项最小实验与强简单基线}
% content block 6:198 p
\bodytext{固定solver、有错初稿和总任务条件，采用文字/\allowbreak{}可视编辑与固定字段澄清/\allowbreak{}学习式澄清两因素比较。强简单基线为字段表单、固定澄清和规则修改器；任务顺序平衡。先用离线回放验接口，再用真实参与者测最终意图一致率、关键漏错与人类总时间，满意度、负荷及操作次数作为辅助。}
% content block 6:199 p
\bodytext{每次修改区分纠正原理解、改变任务、更新偏好与表达不确定，保存原句、前后任务、补丁、确认来源和错误见证。用户点击接受不等于全局最优，新增禁区也不自动证明旧路线违反原始要求。}
% content block 6:84 h3
\Needspace{4\baselineskip}\subsubsection{数据条件和保留条件}
% content block 6:200 p
\bodytext{按参与者、根田块、原始意图和表达模板族留出；生成一千句改写不等于一千个独立意图，学生试点不直接外推农户。样本量依据主终点、预期效应和试点方差确定。若字段表单实用等效，就采用简单方案；若仅满意度增加，结论限于体验；若关键误解与修复时间稳定降低，可形成独立交互贡献。}
% content block 6:21 source
\sourcetext{已有近邻：来源：\href{https://aclanthology.org/2022.findings-acl.28/}{INSPIRED}；\href{https://arxiv.org/abs/2204.05186}{机器人语言纠正}；\href{https://arxiv.org/abs/2401.03701}{ExTraCT}；\href{https://drive.google.com/file/d/1dsLI6ydfXI\_AL8505bH9x8bBM\_JP9zaM/view}{既有视觉VRP研究}。}
% content block 7:7 h1
\clearpage
\hypertarget{part_6}{}
\section{六 数据 成本与风险}
% content block 7:33 h2
\Needspace{6\baselineskip}\subsection{1  数据划分与独立证据单位}
% content block 7:201 p
\bodytext{必须同时区分母实例、求解随机性与算法发现随机性。更多seed可以估计随机求解波动，但不能增加独立田块、独立任务分布或独立发现次数。已被反复查看的BPP确认集、VRPTW测试及NL124继续用于开发和回归；其未来角色要登记，不再默认是未触碰终测。}
% content block 7:85 h3
\Needspace{4\baselineskip}\subsubsection{内层用于学习与选择}
% content block 7:202 p
\bodytext{发现数据用于生成与反馈，推荐训练数据用于拟合选择器，入库验证用于候选保留与信用比较。数据有限时可在开发侧用分组交叉拟合，但特征标准化、特征筛选、调参、阈值与库规模选择都必须留在内层。生成器可以看到训练侧反馈，不能看到外层测试标签。}
% content block 7:86 h3
\Needspace{4\baselineskip}\subsubsection{外层用于评价整个方法}
% content block 7:203 p
\bodytext{如果研究的是完整发现方法，外层每个重复应重新运行相应发现与选择器学习步骤；若只评价一次独立发现后冻结的库，就明确该范围。把同一库的多个solver seed当独立发现，或在外折结果指导生成后仍称它外层测试，都会夸大证据。}
% content block 7:87 h3
\Needspace{4\baselineskip}\subsubsection{怎样形成母组}
% content block 7:8 bullets
\begin{itemize}
\item 同一田块的旋转、点位再采样、容量/\allowbreak{}半径变化、作业行配置和所有seed归同组。
\item 同一母图跨点、行、面任务仍整体留出，才能声称未见场景。
\item BPP同多重集不同顺序相关；同前缀异未来风险组按构造组分析。
\item 按来源、空间邻近与生成程序检查隐性重复；哈希不相同并不自动独立。
\end{itemize}
% content block 7:88 h3
\Needspace{4\baselineskip}\subsubsection{信息可得时点与统计}
% content block 7:204 p
\bodytext{推荐特征只使用决策时可得信息。probe允许获取新信息，但必须支付真实费用。分析采用母组配对差异和适当的不确定性估计；发现机制另考虑独立重启波动。实践意义阈值由农业价值与开发侧变异决定，不在确认集上选一个最有利阈值。}
% content block 7:205 p
\bodytext{分别报告域内新母实例、规模或预算外推、来源与结构分布变化。未显著不等于等效，零观测失败不等于安全认证。样本量与重复由方差、实际效应和主张决定，没有通用的顶刊最低数量。}
% content block 7:34 h2
\Needspace{6\baselineskip}\subsection{2  成本 失败与部署风险}
% content block 7:89 h3
\Needspace{4\baselineskip}\subsubsection{离线与部署两本账}
% content block 7:206 p
\bodytext{离线账包括数据和性能标签、生成调用、修复与失败、准入、候选评测、配置搜索、选择器训练和调参、信用估计，以及独立发现重启。历史父代和预训练资产同时说明完整成本与本轮增量；未知费用保持unknown。}
% content block 7:207 p
\bodytext{部署账覆盖构图或几何编译、特征、probe、选择、求解、后备和checker。若目标是期限内返回可信解，验算也必须在期限内。冷启动与缓存、墙钟与核秒分列；端到端时延还包括排队、通信和结果回收，不能只拿solver内部计时代表用户等待。}
% content block 7:208 p
\bodytext{不要求token、请求数、CPU和墙钟同时机械相等。预先指定与研究问题相符的主资源上限，披露其他实耗及预算—效果曲线。同调用数不等于成本公平；免费并行也不等于同硬件预算。}
% content block 7:90 h3
\Needspace{4\baselineskip}\subsubsection{失败和缺失各有语义}
% content block 7:15 table
\begingroup
\fontsize{9.6}{14.6}\selectfont
\setlength{\tabcolsep}{4pt}
\renewcommand{\arraystretch}{1.18}
\begin{longtable}{@{}>{\raggedright\arraybackslash}p{\dimexpr 0.197406\linewidth-4.7378pt\relax}>{\raggedright\arraybackslash}p{\dimexpr 0.427954\linewidth-10.2709pt\relax}>{\raggedright\arraybackslash}p{\dimexpr 0.374640\linewidth-8.9914pt\relax}@{}}
\toprule
{\sffamily\bfseries 状态} & {\sffamily\bfseries 主分析处理} & {\sffamily\bfseries 不能怎样处理} \\
\midrule
\endfirsthead
\toprule
{\sffamily\bfseries 状态} & {\sffamily\bfseries 主分析处理} & {\sffamily\bfseries 不能怎样处理} \\
\midrule
\endhead
\midrule
\multicolumn{3}{r}{\fontsize{8}{10}\selectfont\color{Muted}续下页}\\
\endfoot
\bottomrule
\endlastfoot
合法截止解 & 按截止时经checker确认的质量计分，保留运行状态。 & 用截止后更好解回填。 \\
\addlinespace[4pt]
无合法解或非法 & 按预先规则计入全分母；另报失败类型。 & 只删除失败后平均质量。 \\
\addlinespace[4pt]
TIMEOUT 或 UNKNOWN & 区分有无合法incumbent、设施故障与状态缺失。 & 等同数学不可行或填零。 \\
\addlinespace[4pt]
尚未评测或日志丢失 & 登记缺失原因；训练侧补测或冻结缺失协议。 & 暗中排最差或按方法挑不同子集。 \\
\addlinespace[4pt]
用户取消或外部修改 & 记录外部中断与新任务版本。 & 作为算法本身失败训练标签。 \\
\end{longtable}
\endgroup
% content block 7:209 p
\bodytext{生成成功率以计划机会为分母，合法候选质量是另一个条件指标。随机部署、多个seed均值和多启动取最好是不同策略；风险校准必须对准实际部署单位。原CRC小样本全回退保留，任何新协议独立冻结。}
% content block 7:91 h3
\Needspace{4\baselineskip}\subsubsection{复用与后训练的摊销}
% content block 7:210 p
\bodytext{离线成本、部署成本之外另报跨任务复用摊销。若一次训练成本为C\_\allowbreak{}train，每个未来设计任务相对冻结强基线净节约ΔC且为正，仅该成本维度的回本需N大于C\_\allowbreak{}train/\allowbreak{}ΔC。一个生成程序服务一万次规划，不等于一万次LLM设计节约；质量和风险须可比，时间、货币、能量分别计算。}
% content block 7:35 h2
\Needspace{6\baselineskip}\subsection{3  风险校准与学习数据资格}
% content block 7:211 p
\bodytext{风险校准用于有清楚风险对象、独立单位和可辨别样本量的问题，不作为每个推荐器的必需模块。CVRP历史432次来自24实例×3算法×2预算×3seed；四折各12训练、6校准、6测试。主设定α=0.1全回退具有结构原因，不能据此认定所有选择学习或后训练无效。}
% content block 7:92 h3
\Needspace{4\baselineskip}\subsubsection{CRC先对齐部署对象}
% content block 7:212 p
\bodytext{已核实现按3个solver seed平均成本判断：超过SBS的1.01倍算有害，且仅在选择非SBS、score超过阈值时计损失。它控制全体实例上的有害切换频率，并非切换后的条件风险或路线regret上界。有限样本修正为(累计损失+1)/\allowbreak{}(n+1)；n=6时下界1/\allowbreak{}7约0.143，高于0.1，因而有限阈值不能通过。n\_\allowbreak{}cal至少9只消除此下界，不保证通过；增加solver seed不增加独立校准实例。}
% content block 7:213 p
\bodytext{单次随机部署、三seed均值和多启动取最好具有不同风险与预算。新协议应先固定部署对象，再核拟合/\allowbreak{}校准隔离、可交换性、阈值损失单调性和端点条件。原CRC为相应条件下的边际期望风险控制，不能改称已选阈值的高概率保证。不得事后放宽阈值只展示有切换的结果。}
% content block 7:16 table
\begingroup
\fontsize{9.6}{14.6}\selectfont
\setlength{\tabcolsep}{4pt}
\renewcommand{\arraystretch}{1.18}
\begin{longtable}{@{}>{\raggedright\arraybackslash}p{\dimexpr 0.187320\linewidth-4.4957pt\relax}>{\raggedright\arraybackslash}p{\dimexpr 0.439481\linewidth-10.5476pt\relax}>{\raggedright\arraybackslash}p{\dimexpr 0.373199\linewidth-8.9568pt\relax}@{}}
\toprule
{\sffamily\bfseries 学习对象} & {\sffamily\bfseries 最小记录单元} & {\sffamily\bfseries 现有缺口} \\
\midrule
\endfirsthead
\toprule
{\sffamily\bfseries 学习对象} & {\sffamily\bfseries 最小记录单元} & {\sffamily\bfseries 现有缺口} \\
\midrule
\endhead
\midrule
\multicolumn{3}{r}{\fontsize{8}{10}\selectfont\color{Muted}续下页}\\
\endfoot
\bottomrule
\endlastfoot
实例选择 & 任务、可用特征、同预算算法结果 & 强库互补、特征时点和干净确认 \\
\addlinespace[4pt]
构造或改进策略 & 状态、合法动作、动作、后继、奖励、终点 & 终局路线不能恢复原solver搜索动作 \\
\addlinespace[4pt]
程序编辑 & 合同、父代、补丁、见证、收益、谱系 & 独立家系和偏好稳定性 \\
\addlinespace[4pt]
语义与交互 & 原意、任务版本、修改、确认和错误见证 & 人工金标、真实参与者及表达族留出 \\
\addlinespace[4pt]
世界模型 & 状态、控制、环境变化、后继和不确定性 & 动作同步、行为覆盖与独立校准 \\
\end{longtable}
\endgroup
% content block 7:22 source
\sourcetext{沿用《研究路线评议》第14、23节。原始CRC：\href{https://arxiv.org/abs/2208.02814}{https:/\allowbreak{}/\allowbreak{}arxiv.org/\allowbreak{}abs/\allowbreak{}2208.02814}；ReCIRC：\href{https://arxiv.org/html/2609.38112v1}{https:/\allowbreak{}/\allowbreak{}arxiv.org/\allowbreak{}html/\allowbreak{}2609.38112v1}。公式是对已有实现的解释，本次未重跑。}
% content block 8:8 h1
\clearpage
\hypertarget{part_7}{}
\section{七 分阶段实施与交付}
% content block 8:36 h2
\Needspace{6\baselineskip}\subsection{1  阶段零与一 任务和运行资格}
% content block 8:93 h3
\Needspace{4\baselineskip}\subsubsection{阶段零 接续已有结果并冻结范围}
% content block 8:214 p
\bodytext{输入是申报原文、最新研究意图、已有代码和原始结果。先完成“需求—升级实现—验收”清单，明确四类农业能力和五项成果；同时建立旧包身份、母组、候选谱系与数据暴露登记。报告、CSV、JSONL和源码哈希能互相对应后，再决定哪些结果可进入新论文。}
% content block 8:215 p
\bodytext{交付物为要求映射、资产清单、缺证清单和冻结协议初版。只有某批记录无法对应时才暂停该批主张，不因统一总账未完就否认其他已有成果。历史成本缺失不补造，新实验开始全量记账。}
% content block 8:94 h3
\Needspace{4\baselineskip}\subsubsection{阶段一 贯通内核并取得合同资格}
% content block 8:216 p
\bodytext{先用强固定点服务或果园方法贯通冻结输入、结果包、独立检查和回放。该合同一旦通过，立即接在线LLM理解与推荐，不等待完整库或行补给模型。行服务、方向状态、资源补给及覆盖扫掠并行适配，分别验收，共用注册、任务身份与运行记录。}
% content block 8:9 bullets
\begin{itemize}
\item 小规模精确或穷举例与独立checker交叉验证；测试有解、无解、非法、超时、取消和恢复。
\item 加入模式链不可拼接、单行超容量、回站不重置总时长、关闭作业穿越不计覆盖等反例。
\item 检查最多K与恰K、任务经过与服务、车体碰撞与工具覆盖等易混语义。
\item 固定当前代码、依赖和检查器版本；未经今日重跑的资产继续标历史验证。
\end{itemize}
% content block 8:95 h3
\Needspace{4\baselineskip}\subsubsection{阶段门和并行工作}
% content block 8:217 p
\bodytext{阶段门按合同滚动检查：当前合同语义、目标与合法性可独立核对后，即可开展它的在线功能、标签和生成反馈。完整性能矩阵只约束相关统计比较，不是所有系统功能的前置。失败先修对应合同与适配器，其他已合格合同继续推进。}
% content block 8:218 p
\bodytext{可同时整理报告与申请素材、准备强基线和界面接线。地图大重构、分布式平台或新模型训练都不是启动科学比较的前置条件。}
% content block 8:37 h2
\Needspace{6\baselineskip}\subsection{2  阶段二与三 强库 在线LLM与真实生成}
% content block 8:96 h3
\Needspace{4\baselineskip}\subsubsection{阶段二 强库和实际推荐空间}
% content block 8:219 p
\bodytext{合同通过资格检查后，登记强传统方法、配置与简明参照，用共同runner生成实例—算法—预算矩阵。能力不合格成员事先门控，不能把排除不适用算法算作学习收益。冻结母组划分，保留独立的新场景。 方法与版本固定：当前PyVRP的ILS实现不能直接称为早期论文的HGS复现。}
% content block 8:220 p
\bodytext{比较训练侧SBS、同预算固定组合、简单规则和廉价监督选择器，另报VBS。分析差异来自终值、早期速度、合法率还是具体农业结构；所有特征、probe和后备费用计入。若固定强法已饱和，保留为默认，转向事先有作业依据的规模、时限或资源条件，而不删除强法。}
% content block 8:97 h3
\Needspace{4\baselineskip}\subsubsection{阶段三 接真实LLM并完成四格归因}
% content block 8:221 p
\bodytext{冻结一个生成模型版本、一个可编辑接口、父代权限、反馈格式和主要资源上限。复用已有真实生成收据、候选谱系和受约束恢复流程，补当前适配器到农业评估的实际接线。先证明真实请求产生的合格源码确实被目标任务执行，而非只运行离线夹具。}
% content block 8:222 p
\bodytext{在线链可从阶段一的首个合格合同开始，独立于离线扩库进度：完成语言理解、缺项处理、能力门控、候选编排与真实求解。离线候选准入后再注册，在线只调用已有ID。两条链按合同汇合，分别评价功能与性能。}
% content block 8:223 p
\bodytext{执行四格及Lcfg、Lret挑战，保留全部计划槽、生成失败、静态拒绝、超时、非法和缺失。若LLM仅产生有效但较弱程序，仍完成生成研究与课题原型，默认部署强传统成员。没有合格生成程序时，只能标生成链有记录，不能标生成算法已接入。}
% content block 8:224 p
\bodytext{学习与模型分支按各自最小实验并行准备，可先做数据字典、接口资格和强简单对照；它们不必等待四格阳性才存在，也不因一次固定池阴性全部暂停。正式扩大算力前须说明要解决的具体不足与独立确认条件。}
% content block 8:98 h3
\Needspace{4\baselineskip}\subsubsection{资源与退出规则}
% content block 8:225 p
\bodytext{使用现有已授权模型账户与调用路径，不在本方案中新增购买或切换计费接口。资源上限由开发侧实际测量后冻结；当失败集中于合同、解析或执行基础设施时先修复，不能把无约束重试称为科学搜索。阶段产物同时进入研究报告与可复现包。}
% content block 8:38 h2
\Needspace{6\baselineskip}\subsection{3  阶段四与五 机制验证与成果形成}
% content block 8:99 h3
\Needspace{4\baselineskip}\subsubsection{阶段四 先同池 再完整发现}
% content block 8:226 p
\bodytext{先用固定候选池做低成本的选库机制比较，保持选择器、库规模、信息和费用一致。测实际质量、合法率、误选损失、被选频次、库稳定性、尾部退化和开销。找到具体机制差异后，再做同预算完整发现重启；若提出原池尚无的新能力假设，也须独立登记与有限验证。}
% content block 8:227 p
\bodytext{闭环中只更改研究机制，其他生成模型、父代、允许组件和实例权限不变。每次重新发现与冻结，外层只评价；开发期试过的机制版本和资源也入账。结果转负时停止相应高成本扩展，不回头使用终测继续调到获胜。}
% content block 8:100 h3
\Needspace{4\baselineskip}\subsubsection{阶段五 使主张与证据相称}
% content block 8:228 p
\bodytext{农业论文围绕完整工作目标、物理相关变量、强方法和可部署收益组织。通用方法论文则须让同一机制在有实质差异的合同上重现，已有BPP、VRPTW/\allowbreak{}CVRP可提供框架和反例，旧测试不能自动当新盲测。}
% content block 8:229 p
\bodytext{根据开发侧方差与预定实际价值确定独立母实例和发现重复；报告全分母、成本范围、母组不确定性和失败边界。主比较缺失时先解决可比较口径，不能靠删除缺失让表格完整。拟主张实机可用，就增加相称的执行证据；只有计算验证就明确范围。}
% content block 8:101 h3
\Needspace{4\baselineskip}\subsubsection{依赖与可并行关系}
% content block 8:10 bullets
\begin{itemize}
\item 需求映射、证据索引、报告骨架与申请准备可从阶段零开始并行。
\item 合同资格先于标签、信用和性能结论；强库及矩阵先于四格归因。
\item 在线LLM功能与界面接线可和离线生成准备并行，共用同一任务版本。
\item 同池机制比较先于解释闭环增益；外层冻结先于正式确认。
\end{itemize}
% content block 8:230 p
\bodytext{阶段按照通过证据推进，不为尚未测量的调用延迟、失败率和计算资源编造工期。每个阶段结束更新剩余工作、实际成本与风险，再决定下一步投入。课题已有可交付成果及时收拢，论文探索不能无限推迟正式交付。}
% content block 8:39 h2
\Needspace{6\baselineskip}\subsection{4  论文组织与贡献边界}
% content block 8:231 p
\bodytext{更近的路线是农业技术研究：围绕一项明确作业，证明正确建模、强算法与推荐如何改善完整农业目标。更高难度的通用方法路线还需要已经验证的新机制及跨合同证据。平台更漂亮、文献更多、重复运行更多，都不能直接补足核心贡献。}
% content block 8:102 h3
\Needspace{4\baselineskip}\subsubsection{农业方向需要补的主证据链}
% content block 8:232 p
\bodytext{当前点服务、导航与运动组件是真实资产，但完整行服务、回补和覆盖联合任务尚缺证据。优先完成一条农业主合同，解释方向、地头、资源与站位怎样影响算法专长；再在有来源差异的新母场景验证。若只有半合成条件，结论就限定于这些条件。}
% content block 8:233 p
\bodytext{CEA官方强调技术增量，农业算法选择已有先例；Soil2Cover显示计算研究可围绕压实与路径联合形成明确农业价值。其他覆盖研究又采用真实田块和导航实验。它们的证据结构可借鉴，不能把某篇实例数、最高提升或实机数量变成所有论文的门槛。[R17至R20]}
% content block 8:103 h3
\Needspace{4\baselineskip}\subsubsection{通用方法方向需要更强归因}
% content block 8:234 p
\bodytext{证明优势来自具体信用或反馈机制，而非更强底层搜索、更多评测、免费的父代、测试泄漏或选择器更换。面对Hydra、实际成员贡献、EoH-S与LACE等近邻，先给出准确方法差异，再用固定池与完整发现两层实验验证。至少让所声称的机制跨越实质不同的任务语义，而不只换地图。}
% content block 8:104 h3
\Needspace{4\baselineskip}\subsubsection{按结果选择论文主张}
% content block 8:11 bullets
\begin{itemize}
\item 单个新程序普遍更强：写单算法发现收益，不能自动写互补库成功。
\item 推荐有效而LLM无额外收益：写可部署推荐，保留生成阴性比较。
\item VBS提高但S不提高：写可观测性、成本或失效边界，需要充分机制证据。
\item 新信用改善独立发现与部署：再论证具体方法增量，并完整对照近邻。
\end{itemize}
% content block 8:235 p
\bodytext{投稿时先判断回答的是农业技术、运输系统还是通用算法问题，再查当时的范围与要求。VRP或LLM关键词不天然适合任何期刊。没有可信的录用概率、顶刊完成百分比或保证成功的固定工期。}
% content block 8:40 h2
\Needspace{6\baselineskip}\subsection{5  交付验收与研究增量}
% content block 8:236 p
\bodytext{结题验收针对可交付能力与成果，论文评价针对可支持的新知识。两套证据共用工程与数据，但不得互相替代。下表作为最终交接检查表，实际学校材料要求仍需在提交前核实。}
% content block 8:17 table
\begingroup
\fontsize{9.6}{14.6}\selectfont
\setlength{\tabcolsep}{4pt}
\renewcommand{\arraystretch}{1.18}
\begin{longtable}{@{}>{\raggedright\arraybackslash}p{\dimexpr 0.185879\linewidth-4.4611pt\relax}>{\raggedright\arraybackslash}p{\dimexpr 0.429395\linewidth-10.3055pt\relax}>{\raggedright\arraybackslash}p{\dimexpr 0.384726\linewidth-9.2334pt\relax}@{}}
\toprule
{\sffamily\bfseries 交付} & {\sffamily\bfseries 完成证据} & {\sffamily\bfseries 与论文的关系} \\
\midrule
\endfirsthead
\toprule
{\sffamily\bfseries 交付} & {\sffamily\bfseries 完成证据} & {\sffamily\bfseries 与论文的关系} \\
\midrule
\endhead
\midrule
\multicolumn{3}{r}{\fontsize{8}{10}\selectfont\color{Muted}续下页}\\
\endfoot
\bottomrule
\endlastfoot
算法库 & 版本化注册表、来源许可、能力测试、成功与失败案例、可复现调用。 & 提供强传统与生成候选；数量不等于创新。 \\
\addlinespace[4pt]
大模型在环原型 & 语言任务→约束→候选ID→真实运行→独立检查→实测比较。 & 分别评价理解、编排、性能选择与生成贡献。 \\
\addlinespace[4pt]
农业案例 & 运输、巡检避障、多点、覆盖能力有对应输入、结果和验算。 & 主文聚焦一个科学问题，其余作系统覆盖与外部检查。 \\
\addlinespace[4pt]
自动设计主线 & 真实请求、父代/\allowbreak{}反馈、候选、演化更新、多问题执行与导出。 & 有效生成、单体胜出、互补及机制创新分别判断。 \\
\addlinespace[4pt]
同名研究报告 & 要求映射、方法、全部主比较、失败、局限、原始证据索引。 & 真实阴性和未完成项不隐藏。 \\
\addlinespace[4pt]
软著与专利申请 & 定版材料、作者/\allowbreak{}权属核对、提交凭证及状态。 & 申请不等于授权；技术主张需有依据。 \\
\end{longtable}
\endgroup
% content block 8:105 h3
\Needspace{4\baselineskip}\subsubsection{交接包应能回答的五个问题}
% content block 8:12 bullets
\begin{itemize}
\item 这是什么版本的任务、程序、数据与评测器？
\item 输入来自哪里，哪部分是真实、合成或未知？
\item 实际跑了哪些计划机会，哪些失败或缺失？
\item 推荐和生成各自带来什么，付出了什么成本？
\item 用同样原件与命令能否重放，哪些结论仍待验证？
\end{itemize}
% content block 8:237 p
\bodytext{两项申报目标都是申请，不是保证授权；材料准备、提交、受理及后续状态分别登记。本轮仅制定研究计划。后续整理技术材料、第三方来源与现有技术，由导师及专业人员评估；未提交不记完成，不预判法律可专利性。交接包保留原件、协议、失败与复现入口，静态示例不标实时求解。}
% content block 9:9 h1
\clearpage
\hypertarget{part_8}{}
\section{附录 来源覆盖与参考依据}
% content block 9:41 h2
\Needspace{6\baselineskip}\subsection{1  来源覆盖与使用边界}
% content block 9:238 p
\bodytext{本方案以三份已有定稿为整合来源。项目事实采用最新《URP研究计划》已核口径；技术选型采用《URP技术方案》；各探索的机制、最小实验及复现入口来自《研究路线评议》。本次不联网、不重读论文库、不运行科研实验。下列覆盖是上轮已记录的阅读范围，不是本次新增阅读。}
% content block 9:239 p
\bodytext{规范索引565项，包括549篇论文与16部教材笔记，来自574份来源的别名归并。已全文阅读551/\allowbreak{}565份规范笔记，即535/\allowbreak{}549篇论文笔记和16/\allowbreak{}16份教材笔记；不等于原论文PDF逐篇全文重读或教材整书重读。另摘读549篇比较卡的“方法、限制、项目关联”，已读16主题综述与4专题。原论文仅回查少数决定性页面。}
% content block 9:240 p
\bodytext{仍未全文读完的14篇规范笔记：P274、P297、P318、P380、P447、P448、P451、P455、P472、P487、P496、P521、P527、P535。保留原清单，不由本次整合自动触发补读。}
% content block 9:241 p
\bodytext{覆盖来自分块回执：M01—M04为120/\allowbreak{}134，M05—M08为136/\allowbreak{}136，M09—M12为178/\allowbreak{}178，M13—M16为117/\allowbreak{}117；旧公共登记尚未统一逐ID回填。M01—M04表内117篇，另有P245、P254、P258全文回执，合计120，不再叠加第一阶段25篇。}
% content block 9:242 p
\bodytext{另有324份待核研究正文，完整读15、部分读2、309未完整读；此前52份入口、综述和报告另计。历史归档198个嵌套成员、4个7z、2个rar未枚举，167个gzip单流未解压。因此不称全部研究文本或历史原稿已读完。}
% content block 9:106 h3
\Needspace{4\baselineskip}\subsubsection{已更正的事实继续沿用}
% content block 9:243 p
\bodytext{新20条EOH对公开EoH为4胜4平12负、多30箱；IID严格18/\allowbreak{}20、技术补齐19/\allowbreak{}20，完整20例主比较仍缺。60槽是生成机会，非独立发现；有限固定池、DAG选择及训练偏好各自保留依赖结构。CALM附录特定基线对应疑点不用于优势论证；不外推否定整篇工作。}
% content block 9:244 p
\bodytext{文献、源码、数据在库，历史实验有账本，当前版本运行合格和田间有效是不同证据层级。版本敏感的软件参数与许可在实施时锁定；本文沿用已有核验结果，不把链接存在写成已经安装或复现。}
% content block 9:42 h2
\Needspace{6\baselineskip}\subsection{2  全主题落实索引}
% content block 9:245 p
\bodytext{本表给出原16主题在统一方案中的位置。“采用”表示拟采用原则或合格组件，不代表已安装、复现或证明增益；“条件探索”均保留最小实验与退出标准。}
% content block 9:18 table
\begingroup
\fontsize{9.6}{14.6}\selectfont
\setlength{\tabcolsep}{4pt}
\renewcommand{\arraystretch}{1.18}
\begin{longtable}{@{}>{\raggedright\arraybackslash}p{\dimexpr 0.302594\linewidth-4.8415pt\relax}>{\raggedright\arraybackslash}p{\dimexpr 0.697406\linewidth-11.1585pt\relax}@{}}
\toprule
{\sffamily\bfseries 主题} & {\sffamily\bfseries 方案中的位置与下一步} \\
\midrule
\endfirsthead
\toprule
{\sffamily\bfseries 主题} & {\sffamily\bfseries 方案中的位置与下一步} \\
\midrule
\endhead
\midrule
\multicolumn{2}{r}{\fontsize{8}{10}\selectfont\color{Muted}续下页}\\
\endfoot
\bottomrule
\endlastfoot
M01 自动算法设计 & 冻结生成器、窄接口、真实反馈；生成修复与进化实验。 \\
\addlinespace[4pt]
M02 库生成与推荐 & 四格、Lcfg/\allowbreak{}Lret、同池信用与独立发现。 \\
\addlinespace[4pt]
M03 agent自改进 & 谱系、失败记忆、固定工具下的harness决策；不扩权或自改验算器。 \\
\addlinespace[4pt]
M04 算法选择与动态配置 & 廉价监督选择先行；运行进展有价值时研究预算控制。 \\
\addlinespace[4pt]
M05 校准与拒判 & 对齐部署事件、校准单位及风险覆盖；CRC阴性保留。 \\
\addlinespace[4pt]
M06 预算与反馈学习 & 固定调度/\allowbreak{}停滞规则、bandit、序贯RL按增量比较。 \\
\addlinespace[4pt]
M07 共演化与泛化 & 母组、分布变化、独立发现；生成困难实例只在开发侧。 \\
\addlinespace[4pt]
M08 统计与程序验证 & 独立checker、全分母、全成本、原意金标准。 \\
\addlinespace[4pt]
M09 学习式路由 & NDS/\allowbreak{}RRNCO/\allowbreak{}GFACS等分构造、改进与控制；先过合同资格。 \\
\addlinespace[4pt]
M10 经典路由与搜索 & PyVRP/\allowbreak{}强局搜、模式搜索与精确小例；排除目标错配。 \\
\addlinespace[4pt]
M11 农业作业与覆盖 & 点、行、面、补给分合同；Fields2Cover等做强对照。 \\
\addlinespace[4pt]
M12 运动与执行 & 姿态、足迹、曲率、局部避障；图合法与现场有效分开。 \\
\addlinespace[4pt]
M13 进化与数学基础 & 保留更新、近似、预算和统计；理论主张贴合真实假设。 \\
\addlinespace[4pt]
M14 代码模型与后训练 & 冻结/\allowbreak{}检索/\allowbreak{}追加搜索/\allowbreak{}SFT或DPO公平比较，开放权重先验训练资格。 \\
\addlinespace[4pt]
M15 架构与强化学习 & 路线RL、预算RL、权重RL分立；NAS资产只支持对应复现。 \\
\addlinespace[4pt]
M16 程序合成与运筹建模 & 结构化合同、受限工具、能力门控及反例修复。 \\
\end{longtable}
\endgroup
% content block 9:43 h2
\Needspace{6\baselineskip}\subsection{3  课题与项目原件来源}
% content block 9:246 p
\bodytext{以下保留原项目相对路径和原始来源身份，供已有交接包检索。文件存在不代表新的方法已执行；本方案所有新实验均为建议。}
% content block 9:107 h3
\Needspace{4\baselineskip}\subsubsection{课题要求}
% content block 9:247 p
\bodytext{[R1] materials/\allowbreak{}项目/\allowbreak{}立项申报书\_\allowbreak{}20260315.docx。第三至四节：算法库、LLM场景理解与约束提取、候选推荐、实际调用和农业案例；第五节：仿真为主；第六节：五项成果。}
% content block 9:248 p
\bodytext{[R2] materials/\allowbreak{}学习/\allowbreak{}LLM自动算法设计入门指南.docx及同名PDF。PDF第20至22页，第六章“实践项目建议”；第21页P3、P5、P6。}
% content block 9:249 p
\bodytext{[R3] materials/\allowbreak{}项目/\allowbreak{}实施方案\_\allowbreak{}20260908.pdf。后续实施取舍依据，不替代正式课题的全部能力要求。}
% content block 9:108 h3
\Needspace{4\baselineskip}\subsubsection{当前保存实验与原件核查}
% content block 9:250 p
\bodytext{[R4] 当前URP只读核查日期2026年10月11日。关键实物如下：}
% content block 9:13 bullets
\begin{itemize}
\item data/\allowbreak{}bpp/\allowbreak{}records/\allowbreak{}MODEL.json、REPORT\_\allowbreak{}ZH.md及analysis/\allowbreak{}confirmation\_\allowbreak{}per\_\allowbreak{}stream.csv、risk\_\allowbreak{}per\_\allowbreak{}stream.csv：冻结SBS、确认与风险重算。
\item experiments/\allowbreak{}results/\allowbreak{}装箱演化/\allowbreak{}原132\_\allowbreak{}P2公开评分56至60\_\allowbreak{}最终增量与封账.zip：生成ledger、真实CLI回执、60槽结果；新序列确认.zip内RESULT\_\allowbreak{}SUMMARY.json及CONFIRMATION\_\allowbreak{}ANALYSIS\_\allowbreak{}FINAL\_\allowbreak{}STRICT\_\allowbreak{}AND\_\allowbreak{}TECHNICAL.json：新数据与缺失口径。
\item experiments/\allowbreak{}results/\allowbreak{}路线选择/\allowbreak{}新64几何前瞻选择器.zip：protocol/\allowbreak{}input\_\allowbreak{}manifest.json、raw/\allowbreak{}train\_\allowbreak{}results.jsonl、raw/\allowbreak{}test\_\allowbreak{}results.jsonl；相应完整双基线中文报告。
\item experiments/\allowbreak{}reports/\allowbreak{}农业路径/\allowbreak{}124块NL算法多seed复验.txt及对应原始包：124母实例、1488槽与重复来源。
\item experiments/\allowbreak{}results/\allowbreak{}服务与运动/\allowbreak{}四路实验\_\allowbreak{}P102容量空间.zip、Dubins固定航向排序.zip及turn\_\allowbreak{}modes结果；experiments/\allowbreak{}results/\allowbreak{}组件/\allowbreak{}orchard\_\allowbreak{}navigation/\allowbreak{}results/\allowbreak{}main/\allowbreak{}raw.jsonl、replay.jsonl。
\end{itemize}
% content block 9:251 p
\bodytext{当前工程边界：src/\allowbreak{}research/\allowbreak{}README.md、src/\allowbreak{}backend/\allowbreak{}backend/\allowbreak{}server.py；农业适配限制：experiments/\allowbreak{}results/\allowbreak{}组件/\allowbreak{}vehicle\_\allowbreak{}routing/\allowbreak{}initialization\_\allowbreak{}ablation\_\allowbreak{}v1/\allowbreak{}adapter.py。原件中的历史状态和后续独立审计分别保留。}
% content block 9:252 p
\bodytext{[R21] 文献定位：literature/\allowbreak{}index.csv；literature/\allowbreak{}资料/\allowbreak{}来源合并.json；literature/\allowbreak{}笔记/\allowbreak{}下16主题的主题综述、比较表及逐篇笔记。覆盖依据：2026年10月11日阅读登记与M01至M16分块回执。规范未读以本方案覆盖附录列出的14个ID为准。}
% content block 9:44 h2
\Needspace{6\baselineskip}\subsection{4  方法与工程来源}
% content block 9:253 p
\bodytext{以下为本方案主要方法判断的原始论文或官方工程入口。版本敏感实现应在实际复现时固定提交；参考入口可读取不等于已在URP安装或运行。}
% content block 9:254 p
\bodytext{[R5] Romera-Paredes等，FunSearch，Nature 2024。LLM程序发现、独立评估与未见任务验证。\newline \href{https://doi.org/10.1038/s41586-023-06924-6}{论文}}
% content block 9:255 p
\bodytext{[R6] EoH论文与官方实现；ReEvo，NeurIPS 2024。思想和代码演化、反思式生成反馈。\newline \href{https://arxiv.org/abs/2401.02051}{论文}\newline \href{https://github.com/FeiLiu36/EoH}{源码}\newline \href{https://proceedings.neurips.cc/paper\_files/paper/2024/hash/4ced59d480e07d290b6f29fc8798f195-Abstract-Conference.html}{论文}}
% content block 9:256 p
\bodytext{[R7] Bischl等，ASlib，Artificial Intelligence 2016。统一算法选择场景、性能与特征费用。\newline \href{https://doi.org/10.1016/j.artint.2016.04.003}{论文}\newline \href{https://arxiv.org/html/1506.02465v3}{论文}}
% content block 9:257 p
\bodytext{[R8] AutoFolio、irace与LLM4AD官方工程。分别用于选择、配置和生成评估适配。\newline \href{https://github.com/automl/AutoFolio}{源码}\newline \href{https://mlopez-ibanez.github.io/irace/}{官方来源}\newline \href{https://github.com/Optima-CityU/LLM4AD}{源码}}
% content block 9:258 p
\bodytext{[R9] Xu、Hoos、Leyton-Brown，Hydra，2010；SATzilla成员贡献研究，2012。互补配置与实际选择，成员省略后重新训练。\newline \href{https://www.cs.ubc.ca/~hoos/Publ/XuEtAl10.pdf}{官方来源}\newline \href{https://www.cs.ubc.ca/~kevinlb/papers/2012-SAT-EvaluatingSolverContributions.pdf}{官方来源} \href{https://www.cs.ubc.ca/labs/algorithms/Projects/Hydra/}{Hydra作者资源}；\href{https://www.cs.ubc.ca/labs/algorithms/Projects/SATzilla/}{SATzilla作者资源}。}
% content block 9:259 p
\bodytext{[R10] EoH-S正式论文。互补启发式集合与CPI。\newline \href{https://ojs.aaai.org/index.php/AAAI/article/download/41038/44999}{论文}}
% content block 9:260 p
\bodytext{[R11] LACE，Nature Machine Intelligence 2026及官方源码；本方案采用其合同、组合和rank-MILP近邻定位，具体实现边界以固定源码审计为据。\newline \href{https://www.nature.com/articles/s42256-026-01307-8}{论文}\newline \href{https://github.com/PJ-NTU/LACE}{源码} \href{https://github.com/PJ-NTU/LACE/blob/d772b74690075f641be8077e69748ec482e9febf/lace/evolution/milp\_selection.py}{已核rank-MILP固定实现}。}
% content block 9:261 p
\bodytext{[R12] Fields2Cover与PyVRP官方工程。覆盖规划、车辆路线及补给能力需通过URP窄适配和独立验算。\newline \href{https://github.com/Fields2Cover/Fields2Cover}{源码}\newline \href{https://github.com/PyVRP/PyVRP}{源码}}
% content block 9:45 h2
\Needspace{6\baselineskip}\subsection{5  农业与交互来源}
% content block 9:262 p
\bodytext{[R13] Vidal，NEARP作者稿与行服务参考。方向与服务模式不等同完整农业补给，须独立适配。\newline \href{https://w1.cirrelt.ca/~vidalt/papers/NEARP-Vidal-Revision-WP.pdf}{官方来源}}
% content block 9:263 p
\bodytext{[R14] RL4CO与DACBench官方工程；IQL原论文。环境接口、动态配置及离线学习提供研究入口，不证明当前数据已经适配。\newline \href{https://rl4co.ai4co.org/docs/content/api/envs/base/}{官方来源}\newline \href{https://github.com/automl/DACBench}{源码}\newline \href{https://arxiv.org/abs/2110.06169}{论文}}
% content block 9:264 p
\bodytext{[R15] OptiGuide官方工程，LLM-Modulo原论文，PICARD（EMNLP 2021）。分别借鉴真实反馈、验证权分离和受约束接口。\newline \href{https://github.com/microsoft/OptiGuide/tree/main/what-if}{源码}\newline \href{https://arxiv.org/html/2402.01817v3}{论文}\newline \href{https://aclanthology.org/2021.emnlp-main.779/}{官方来源}}
% content block 9:265 p
\bodytext{[R16] RELLIS-3D与TartanDrive 2.0官方数据入口。\newline \href{https://github.com/unmannedlab/RELLIS-3D}{源码}\newline \href{https://theairlab.org/TartanDrive2/}{官方来源}}
% content block 9:266 p
\bodytext{[R17] CEA 2019农业路径与算法选择相关论文。农业加算法选择已有先例。\newline \href{https://www.sciencedirect.com/science/article/abs/pii/S0168169918303867}{官方来源}}
% content block 9:267 p
\bodytext{[R18] JRM 2026农业SDVRP统一框架及算法/\allowbreak{}参数选择，DOI 10.20965/\allowbreak{}jrm.2026.p0413。\newline \href{https://www.jstage.jst.go.jp/article/jrobomech/38/2/38\_413/\_pdf}{官方来源}}
% content block 9:268 p
\bodytext{[R19] Mier等，Soil2Cover，Precision Agriculture 2025。1000田块计算验证；最高改善不等于全体平均，不据其样本数规定URP门槛。\newline \href{https://link.springer.com/article/10.1007/s11119-025-10250-4}{官方来源}}
% content block 9:269 p
\bodytext{[R20] CEA官方范围及2025覆盖规划论文；Precision Agriculture范围；TR-C官方范围。期刊适配按研究问题与实际证据判断。\newline \href{https://shop.elsevier.com/journals/computers-and-electronics-in-agriculture/0168-1699}{官方来源}\newline \href{https://www.sciencedirect.com/science/article/pii/S0168169925000523}{官方来源}\newline \href{https://link.springer.com/journal/11119/aims-and-scope}{官方来源}\newline \href{https://shop.elsevier.com/journals/transportation-research-part-c-emerging-technologies/0968-090X}{官方来源}}
% content block 9:270 p
\bodytext{正文中的原始结果数字由当前文件核查支持；期刊样例用于比较证据结构，不是通用录用门槛。后续正式论文应依据最终采用的方法补齐完整书目信息、固定版本、数据与代码可用性说明。}
% content block 9:46 h2
\Needspace{6\baselineskip}\subsection{6  既有原文缺口与条件性应用}
% content block 9:271 p
\bodytext{沿用《研究路线评议》已核清单，按实际采用的研究分支安排补核；不在本次整合继续调查。原文缺口不等于已有方法或源码完全缺失。}
% content block 9:19 table
\begingroup
\fontsize{9.6}{14.6}\selectfont
\setlength{\tabcolsep}{4pt}
\renewcommand{\arraystretch}{1.18}
\begin{longtable}{@{}>{\raggedright\arraybackslash}p{\dimexpr 0.381844\linewidth-9.1643pt\relax}>{\raggedright\arraybackslash}p{\dimexpr 0.273775\linewidth-6.5706pt\relax}>{\raggedright\arraybackslash}p{\dimexpr 0.344380\linewidth-8.2651pt\relax}@{}}
\toprule
{\sffamily\bfseries 来源} & {\sffamily\bfseries 原记录状态} & {\sffamily\bfseries 何时需要} \\
\midrule
\endfirsthead
\toprule
{\sffamily\bfseries 来源} & {\sffamily\bfseries 原记录状态} & {\sffamily\bfseries 何时需要} \\
\midrule
\endhead
\midrule
\multicolumn{3}{r}{\fontsize{8}{10}\selectfont\color{Muted}续下页}\\
\endfoot
\bottomrule
\endlastfoot
Accorsi等 VRP机器学习计算测试指南\newline DOI 10.1016/\allowbreak{}j.orl.2022.01.018 & 未发现独立原文，已有官方与开放入口 & 冻结强OR基线、完整计时和统计协议时 \\
\addlinespace[4pt]
LACE\newline DOI 10.1038/\allowbreak{}s42256-026-01307-8 & 主文未在库；摘要与固定代码已核 & 精确论证互补生成机制与新颖性时 \\
\addlinespace[4pt]
粮食收获路线与指标\newline DOI 10.3390/\allowbreak{}agronomy12051151 & 原库仅引用，已有开放原文入口 & 行作业和资源成本合同确定时 \\
\addlinespace[4pt]
草莓收获辅助系统\newline DOI 10.1002/\allowbreak{}rob.22106 & 未找到匹配原文，有开放全文入口 & 启动动态收获请求或现场指标研究时 \\
\addlinespace[4pt]
喷雾车与补给车同步\newline DOI 10.1016/\allowbreak{}j.cor.2023.106454 & 未找到独立原文，前稿身份不自动合并 & 移动补给与会合分支，不是固定仓库补满前置 \\
\addlinespace[4pt]
可拆卸农具协作车队\newline DOI 10.1016/\allowbreak{}j.cor.2025.107252 & 未在库，有开放正式版入口 & 农具交换、合作车队或公平性研究时 \\
\addlinespace[4pt]
农业覆盖评价基准\newline DOI 10.1016/\allowbreak{}j.biosystemseng.2020.08.007 & 有引用未找到独立原文 & 现有Fields2Benchmark协议不足时 \\
\addlinespace[4pt]
DAR DPO\newline arXiv 2507.10614v2 & 有题名引用，未找到独立原文 & 正式使用程序档案开展偏好训练时 \\
\end{longtable}
\endgroup
% content block 9:272 p
\bodytext{P042 EvoTune指定v4、P446农业算法选择、P202 Fields2Benchmark、P203 Fields2Cover、P102车具容量续航、P511服务同步、P474 Soil2Cover均已有原文或相应资产，不重复列为新缺口。农业算法选择、补给、合同验证和行序学习已有直接先例，应转化为对照与验收。}
% content block 9:273 p
\bodytext{对应入口如下，均沿用原稿，不表示本次新增核验。已有EvoTune、农业选择、覆盖、车具和土壤目标来源另列，以便直接转为实验对照。}
% content block 9:23 source
\sourcetext{VRP测试指南官方/\allowbreak{}开放稿：\href{https://publications.polymtl.ca/50935/}{入口1} \href{https://arxiv.org/abs/2109.13983}{入口2}}
% content block 9:24 source
\sourcetext{LACE正式页与源码：\href{https://www.nature.com/articles/s42256-026-01307-8}{入口1} \href{https://github.com/PJ-NTU/LACE}{入口2}}
% content block 9:25 source
\sourcetext{粮食收获开放原文：\href{https://www.mdpi.com/2073-4395/12/5/1151}{入口1}}
% content block 9:26 source
\sourcetext{草莓收获辅助正式页与开放全文：\href{https://onlinelibrary.wiley.com/doi/full/10.1002/rob.22106}{入口1} \href{https://s.gifford.ucdavis.edu/uploads/pub/2023/03/10/vougioukasrobots.pdf}{入口2}}
% content block 9:27 source
\sourcetext{移动补给正式页：\href{https://www.sciencedirect.com/science/article/pii/S0305054823003180}{入口1}}
% content block 9:28 source
\sourcetext{合作农具车队正式页与开放全文：\href{https://www.sciencedirect.com/science/article/pii/S0305054825002813}{入口1} \href{https://burjcdigital.urjc.es/bitstreams/eab902dd-2b44-4f09-acc9-6507f194b8ec/download}{入口2}}
% content block 9:29 source
\sourcetext{覆盖基准题录：\href{https://pure.au.dk/portal/en/publications/method-and-bench-marking-framework-for-coverage-path-planning-in-/}{入口1}}
% content block 9:30 source
\sourcetext{DAR DPO指定版本：\href{https://arxiv.org/abs/2507.10614v2}{入口1}}
% content block 9:31 source
\sourcetext{\textbf{Routing algorithm selection for field coverage planning based on field shape and fleet size}。Seyyedhasani、Dvorak、Roemmele，Computers and Electronics in Agriculture 156，523–529，2019。DOI：10.1016/\allowbreak{}j.compag.2018.12.002。\href{https://scholars.uky.edu/en/publications/routing-algorithm-selection-for-field-coverage-planning-based-on-/}{作者机构记录}。原文和笔记在库；本轮在线核摘要，未另核开放全文。}
% content block 9:32 source
\sourcetext{\textbf{Fields2Benchmark: An Open-Source Benchmark for Coverage Path Planning Methods in Agriculture}，Mier等，2025。DOI：10.1016/\allowbreak{}j.atech.2025.101156。\href{https://edepot.wur.nl/698168}{作者原文}、\href{https://zenodo.org/records/14524735}{田界数据}。P202原文、笔记及350田块wkt.zip已在库；读任务／指标与局限，核和NL124交叠，补覆盖验收。}
% content block 9:33 source
\sourcetext{\textbf{Fields2Cover: An Open-Source Coverage Path Planning Library for Unmanned Agricultural Vehicles}，Mier等，2023。DOI：10.1109/\allowbreak{}LRA.2023.3248439。\href{https://arxiv.org/abs/2210.07838}{开放作者稿}、\href{https://github.com/Fields2Cover/Fields2Cover}{官方工程}。P203原文、笔记和代码已在库；读完整流程，固定车宽／幅宽、地头、行与Dubins版本，统一覆盖指标及整机几何回查。}
% content block 9:34 source
\sourcetext{\textbf{Capacitated Agriculture Fleet Vehicle Routing with Implements and Limited Autonomy: A model and a two-phase solution approach}。López-Sánchez等，ICRA 2025。DOI：10.1109/\allowbreak{}ICRA55743.2025.11127753。\href{https://burjcdigital.urjc.es/items/4bfae71f-638f-4c78-80e5-79b8274f5a39}{作者机构入口}。HAL稿题名、作者和会议吻合，归档与主库哈希一致，稿内未印出版DOI。}
% content block 9:35 source
\sourcetext{\textbf{The Vehicle Routing Problem in Field Logistics Part II}，Bochtis、Sørensen，2010。DOI：10.1016/\allowbreak{}j.biosystemseng.2009.10.006。\href{https://doi.org/10.1016/j.biosystemseng.2009.10.006}{出版入口}。原文与笔记在库，开放全文未另核；读支持车辆时空请求，落实释放、位置、服务与等待账本，未确认对应项目实验。}
% content block 9:36 source
\sourcetext{\textbf{Soil2Cover: Coverage path planning minimizing soil compaction for sustainable agriculture}，Mier等，Precision Agriculture 26，57，2025。DOI：10.1007/\allowbreak{}s11119-025-10250-4。\href{https://doi.org/10.1007/s11119-025-10250-4}{出版社入口}。原文与笔记在库；选土壤目标时读压实模型、通行次数与评价，区分代理优化和实测影响。未确认独立参考仓或项目压实优化实验。}
% content block 9:47 h2
\Needspace{6\baselineskip}\subsection{7  技术实现来源}
% content block 9:274 p
\bodytext{以下T编号对应技术章节；官方入口及版本判断沿用原技术方案截至2026年10月11日的核验。实施时锁定实际版本与许可。}
% content block 9:1 reference
\referencetext{[T1] OpenLayers 编辑能力：\href{https://openlayers.org/en/latest/apidoc/module-ol\_interaction\_Draw-Draw.html}{Draw}、\href{https://openlayers.org/en/latest/apidoc/module-ol\_interaction\_Modify-Modify.html}{Modify}、\href{https://openlayers.org/en/latest/apidoc/module-ol\_interaction\_Snap-Snap.html}{Snap}。\href{https://openlayers.org/en/latest/examples/webgl-vector-layer.html}{WebGL Vector 示例的非稳定 API 提示}。}
% content block 9:2 reference
\referencetext{[T2] OpenLayers 坐标与栅格：\href{https://openlayers.org/en/latest/apidoc/module-ol\_proj\_Projection-Projection.html}{Projection 与内建 UTM 精度警告}、\href{https://openlayers.org/en/latest/apidoc/module-ol\_proj\_proj4.html}{proj4 注册}、\href{https://openlayers.org/en/latest/examples/static-image.html}{自定义平面示例}、\href{https://openlayers.org/en/latest/apidoc/module-ol\_source\_GeoTIFF-GeoTIFFSource.html}{GeoTIFF Source}。}
% content block 9:3 reference
\referencetext{[T3] MapLibre：\href{https://maplibre.org/maplibre-style-spec/types/\#projectiondefinition}{投影定义}、\href{https://maplibre.org/maplibre-gl-js/docs/guides/large-data/}{大 GeoJSON 指南}、\href{https://maplibre.org/maplibre-gl-js/docs/guides/v5-to-v6-migration-guide/}{v6 迁移指南}、\href{https://github.com/geomatico/maplibre-cog-protocol}{常用 COG 插件的 EPSG 3857 限制}。}
% content block 9:4 reference
\referencetext{[T4] Terra Draw：\href{https://github.com/JamesLMilner/terra-draw/blob/main/guides/4.MODES.md}{模式与验证}、\href{https://github.com/JamesLMilner/terra-draw/blob/main/guides/6.EVENTS.md}{两级历史与事件}、\href{https://raw.githubusercontent.com/JamesLMilner/terra-draw/main/packages/terra-draw-maplibre-gl-adapter/package.json}{adapter 包依赖}、\href{https://raw.githubusercontent.com/watergis/maplibre-gl-terradraw/main/package.json}{工具栏包依赖}。}
% content block 9:5 reference
\referencetext{[T5] Geoman 与 Leaflet：\href{https://geoman.io/docs/maplibre}{MapLibre Geoman 的 Free 和 Pro}、\href{https://geoman.io/docs/leaflet/modes/edit-mode}{Leaflet 编辑模式}、\href{https://leafletjs.com/examples/crs-simple/}{Leaflet 本地平面教程}。}
% content block 9:6 reference
\referencetext{[T6] deck.gl：\href{https://deck.gl/docs/developer-guide/coordinate-systems}{坐标系统}、\href{https://deck.gl/docs/developer-guide/performance}{性能指南}、\href{https://visgl.github.io/deck.gl-community/docs/modules/editable-layers}{社区 editable-layers 维护说明}。}
% content block 9:7 reference
\referencetext{[T7] 三维与完整 GIS：\href{https://cesium.com/platform/cesiumjs/}{CesiumJS}、\href{https://docs.mapstore.geosolutionsgroup.com/en/latest/developer-guide/infrastructure-and-general-architecture/}{MapStore 架构}。}
% content block 9:8 reference
\referencetext{[T8] 计算几何与坐标：\href{https://pyproj4.github.io/pyproj/stable/api/transformer.html}{pyproj Transformer}、\href{https://proj.org/en/stable/operations/projections/utm.html}{PROJ UTM}、\href{https://proj.org/en/stable/usage/network.html}{PROJ 离线与网络格网}、\href{https://shapely.readthedocs.io/en/stable/manual.html}{Shapely 手册}。}
% content block 9:9 reference
\referencetext{[T9] 标准和原始基准：\href{https://datatracker.ietf.org/doc/html/rfc7946}{RFC 7946 GeoJSON}、\href{https://zenodo.org/records/14524735}{Fields2Benchmark 发布页及各国来源许可}。}
% content block 9:10 reference
\referencetext{[T10] 文件与栅格：\href{https://www.geopackage.org/spec140/}{GeoPackage 1.4 标准}、\href{https://gdal.org/en/stable/drivers/raster/cog.html}{GDAL COG}、\href{https://rasterio.readthedocs.io/en/latest/topics/windowed-rw.html}{Rasterio 按窗读取}。}
% content block 9:11 reference
\referencetext{[T11] 离线底图：\href{https://docs.protomaps.com/pmtiles/}{PMTiles}、\href{https://operations.osmfoundation.org/policies/tiles/}{OSM 栅格政策}、\href{https://operations.osmfoundation.org/policies/vector/}{OSM 矢量政策}。}
% content block 9:12 reference
\referencetext{[T12] 采集与数据检查：\href{https://docs.qgis.org/3.44/en/docs/user\_manual/managing\_data\_source/georeferencer.html}{QGIS Georeferencer}、\href{https://github.com/opengisch/QField}{QField}、\href{https://developer.mozilla.org/en-US/docs/Web/API/Storage\_API/Storage\_quotas\_and\_eviction\_criteria}{浏览器存储配额与清理}。}
% content block 9:13 reference
\referencetext{[T13] 覆盖与点服务：\href{https://github.com/Fields2Cover/Fields2Cover}{Fields2Cover}、\href{https://fields2cover.github.io/source/tutorials/basic\_types.html}{Route 与 Path 类型}、\href{https://github.com/PyVRP/PyVRP/releases}{PyVRP 发布记录}。}
% content block 9:14 reference
\referencetext{[T14] 存储与桌面：\href{https://sqlite.org/whentouse.html}{SQLite 适用范围}、\href{https://sqlite.org/wal.html}{WAL}、\href{https://v2.tauri.app/reference/webview-versions/}{Tauri WebView}、\href{https://v2.tauri.app/develop/sidecar/}{Tauri Sidecar}、\href{https://www.electronjs.org/docs/latest/tutorial/security}{Electron 安全清单}。}
% content block 9:15 reference
\referencetext{[T15] 外部观测器：\href{https://rerun.io/docs/howto/visualization/geospatial-data}{Rerun 地理数据能力}、\href{https://foxglove.dev/blog/foxglove-2-0-unifying-robotics-observability}{Foxglove 2024 产品与开源调整}。}
% content block 9:16 reference
\referencetext{[T16] 当前产品与存档：src/\allowbreak{}backend/\allowbreak{}backend/\allowbreak{}server.py、src/\allowbreak{}research/\allowbreak{}README.md；experiments/\allowbreak{}provenance/\allowbreak{}记录/\allowbreak{}综合/\allowbreak{}URP\_\allowbreak{}实验项目.zip 内 URP/\allowbreak{}web/\allowbreak{}planner.html 与 planner.js。10月11日只读核查确认历史原型与当前React未统一接通。}
% content block 9:17 reference
\referencetext{[T17] 行与区域任务：\href{https://github.com/open-navigation/opennav\_coverage}{Open Navigation Coverage}、\href{https://w1.cirrelt.ca/~vidalt/papers/NEARP-Vidal-Revision-WP.pdf}{Vidal 等 NEARP 服务模式论文}。}
% content block 9:18 reference
\referencetext{[T18] 规划约束：\href{https://developers.google.com/optimization/routing/dimensions?hl=en}{OR-Tools Dimensions}、\href{https://ompl.kavrakilab.org/stateValidation.html}{OMPL 状态与运动有效性}。}
% content block 9:19 reference
\referencetext{[T19] API 与进程：\href{https://fastapi.tiangolo.com/tutorial/background-tasks/}{FastAPI 后台任务边界}、\href{https://docs.pydantic.dev/latest/concepts/models/}{Pydantic 模型}、\href{https://docs.python.org/3/library/subprocess.html}{Python subprocess}。}
% content block 9:20 reference
\referencetext{[T20] 执行隔离与资源：\href{https://docs.docker.com/engine/security/}{Docker 安全}、\href{https://docs.ray.io/en/latest/ray-core/scheduling/resources.html}{Ray 逻辑资源}。}
% content block 9:21 reference
\referencetext{[T21] 队列和分析：\href{https://docs.celeryq.dev/en/stable/userguide/tasks.html}{Celery 任务}、\href{https://docs.celeryq.dev/en/stable/userguide/workers.html}{worker 与取消}、\href{https://docs.dask.org/en/stable/best-practices.html}{Dask 最佳实践}、\href{https://sqlite.org/releaselog/3\_51\_3.html}{SQLite WAL 修复发布记录}。}
% content block 9:22 reference
\referencetext{[T26] 联动视图：\href{https://reactflow.dev/index}{React Flow}、\href{https://echarts.apache.org/handbook/en/concepts/data-transform/}{ECharts 数据变换}、\href{https://github.com/vega/vega-lite/blob/main/site/docs/parameter/select.md}{Vega-Lite 选择交互}。已有图表可继续使用，无须为联动而换库。}
% content block 9:23 reference
\referencetext{[T28] 农业原始证据：experiments/\allowbreak{}reports/\allowbreak{}农业路径/\allowbreak{}124块NL算法多seed复验.txt；experiments/\allowbreak{}results/\allowbreak{}服务与运动/\allowbreak{}四路实验\_\allowbreak{}P102容量空间.zip、Dubins固定航向排序.zip；src/\allowbreak{}research/\allowbreak{}src/\allowbreak{}urp\_\allowbreak{}research/\allowbreak{}turn\_\allowbreak{}modes/\allowbreak{}checker.py；experiments/\allowbreak{}results/\allowbreak{}组件/\allowbreak{}orchard\_\allowbreak{}navigation/\allowbreak{}results/\allowbreak{}main/\allowbreak{}；所查回补适配器为 experiments/\allowbreak{}results/\allowbreak{}组件/\allowbreak{}vehicle\_\allowbreak{}routing/\allowbreak{}initialization\_\allowbreak{}ablation\_\allowbreak{}v1/\allowbreak{}adapter.py。}
% content block 9:24 reference
\referencetext{[Tinteraction1] 交互修改合同补充依据既有《研究路线评议》第11与20节；原文链接：\href{https://aclanthology.org/2022.findings-acl.28/}{INSPIRED}、\href{https://arxiv.org/abs/2204.05186}{机器人语言纠正}、\href{https://arxiv.org/abs/2401.03701}{ExTraCT}。}
% content block 9:25 reference
\referencetext{[Tinteraction2] 逐来源坐标转换示例：\href{https://github.com/ucmercedrobotics/gpt-mission-planner/blob/6c0aa5706f58fb22bd98b6ddd1968be1681c1d16/app/resources/context/wheeled\_bots/reza\_medium.geojson}{As You Wish地图}与\href{https://github.com/ucmercedrobotics/gpt-mission-planner/blob/6c0aa5706f58fb22bd98b6ddd1968be1681c1d16/app/gpt\_outputs/gpt\_example.xml\#L11-L44}{同提交XML}。}
% content block 9:48 h2
\Needspace{6\baselineskip}\subsection{8  探索方法来源}
% content block 9:275 p
\bodytext{ER编号对应探索章节，沿用既有评议来源和已核前次NDS入口。本次未重新核验链接或软件版本。}
% content block 9:26 reference
\referencetext{[ER01] \href{https://arxiv.org/html/2503.16159v1}{RRNCO}。}
% content block 9:27 reference
\referencetext{[ER02] \href{https://arxiv.org/html/2509.10303v3}{CDQAC v3}。}
% content block 9:28 reference
\referencetext{[ER03] \href{https://github.com/jesserem/CDQAC\_scheduling/blob/aca0896aae7b715606b8f47e870e4489e50ebf5c/cdqac/methods/iql.py}{CDQAC离散IQL固定源码}。}
% content block 9:29 reference
\referencetext{[ER04] \href{https://github.com/jesserem/CDQAC\_scheduling/blob/aca0896aae7b715606b8f47e870e4489e50ebf5c/cdqac/buffer/buffer.py\#L138-L173}{CDQAC轨迹buffer固定源码}。}
% content block 9:30 reference
\referencetext{[ER05] \href{https://arxiv.org/abs/2211.00759}{DR-ALNS}。}
% content block 9:31 reference
\referencetext{[ER06] \href{https://papers.neurips.cc/paper\_files/paper/2021/hash/dc9fa5f217a1e57b8a6adeb065560b38-Abstract.html}{Learning to Delegate}。}
% content block 9:32 reference
\referencetext{[ER07] \href{https://arxiv.org/html/2608.03636v1}{MuEvo}。}
% content block 9:33 reference
\referencetext{[ER08] \href{https://arxiv.org/html/2603.08493v1}{PolarBear}。}
% content block 9:34 reference
\referencetext{[ER09] \href{https://proceedings.mlr.press/v235/liu24bs.html}{EoH}。}
% content block 9:35 reference
\referencetext{[ER10] \href{https://arxiv.org/html/2609.15820v2}{AlgoEvo v2}。}
% content block 9:36 reference
\referencetext{[ER11] \href{https://arxiv.org/html/2605.08756v2}{AHD Agent v2}。}
% content block 9:37 reference
\referencetext{[ER12] \href{https://github.com/large-loris-models/chopchop/blob/681083a6fd921ac9cbaf984db628cf92eb019a3f/llm/run\_llm.py}{ChopChop约束与计时固定源码}。}
% content block 9:38 reference
\referencetext{[ER13] \href{https://github.com/pmorvalho/LLM-CEGIS-Repair/blob/7fe749b16a7a1367f2a519492a5d8689af06a896/repair.py}{LLM-CEGIS-Repair修复循环}。}
% content block 9:39 reference
\referencetext{[ER14] \href{https://github.com/pmorvalho/LLM-CEGIS-Repair/blob/7fe749b16a7a1367f2a519492a5d8689af06a896/mentor/interpreter.py}{LLM-CEGIS-Repair Python判定}。}
% content block 9:40 reference
\referencetext{[ER15] \href{https://github.com/vidalt/HGS-CARP/blob/08fde6a5672c3814dca1b4450aa9e5968b1d273a/Program/Params.cpp\#L1041-L1097}{HGS-CARP成本与转向固定源码}。}
% content block 9:41 reference
\referencetext{[ER16] \href{https://github.com/vidalt/HGS-CARP/blob/08fde6a5672c3814dca1b4450aa9e5968b1d273a/Program/SeqData.cpp\#L514-L538}{HGS-CARP子序列mode表}。}
% content block 9:42 reference
\referencetext{[ER17] \href{https://arxiv.org/html/2504.05108v4}{EvoTune v4}。}
% content block 9:43 reference
\referencetext{[ER18] \href{https://github.com/CLAIRE-Labo/EvoTune/blob/683b245251e96332892a52d2ad57e262f9d6e4c2/src/experiments/main.py}{EvoTune固定入口}。}
% content block 9:44 reference
\referencetext{[ER19] \href{https://github.com/CLAIRE-Labo/EvoTune/blob/683b245251e96332892a52d2ad57e262f9d6e4c2/src/packing/train/train.py}{EvoTune固定训练器}。}
% content block 9:45 reference
\referencetext{[ER20] \href{https://proceedings.iclr.cc/paper\_files/paper/2026/hash/75ed738c083f8941cfc990245e4f6265-Abstract-Conference.html}{CALM ICLR 2026正式版}。}
% content block 9:46 reference
\referencetext{[ER21] \href{https://github.com/whxru/CALM/blob/3c40dbd78f2f6439bcc7cc2cbac334ab28a6f8f3/problems/cvrp\_aco.py}{CALM CVRP宿主固定源码}。}
% content block 9:47 reference
\referencetext{[ER22] \href{https://arxiv.org/html/2610.06269v1}{Guidance-TTT}。}
% content block 9:48 reference
\referencetext{[ER23] \href{https://github.com/Human-Agent-Society/reef/blob/ed8b415cf6eb1051902044428d6ac49c9ad0e3d9/recipes/tttd/examples/guidance\_ttt/README.md}{Guidance-TTT固定工程说明}。}
% content block 9:49 reference
\referencetext{[ER24] \href{https://arxiv.org/html/2605.07039v1}{PACEvolve++}。}
% content block 9:50 reference
\referencetext{[ER25] \href{https://arxiv.org/html/2507.10614v2}{DAR-DPO v2}。}
% content block 9:51 reference
\referencetext{[ER26] \href{https://arxiv.org/html/2609.39383v1}{From Search to Signal}。}
% content block 9:52 reference
\referencetext{[ER27] \href{https://arxiv.org/html/2609.38757v1}{PCPO}。}
% content block 9:53 reference
\referencetext{[ER28] \href{https://huggingface.co/Qwen/Qwen3.5-4B}{Qwen3.5-4B官方卡}。}
% content block 9:54 reference
\referencetext{[ER29] \href{https://huggingface.co/google/gemma-4-12B-it}{Gemma 4 12B官方卡}。}
% content block 9:55 reference
\referencetext{[ER30] \href{https://proceedings.mlr.press/v258/kim25a.html}{GFACS}。}
% content block 9:56 reference
\referencetext{[ER31] \href{https://github.com/ai4co/gfacs/blob/39d88cf2a5f9eb81eeea5a43ee57ff74ac8af01e/cvrp\_nls/train.py}{GFACS固定训练代码}。}
% content block 9:57 reference
\referencetext{[ER32] \href{https://github.com/ai4co/gfacs/blob/39d88cf2a5f9eb81eeea5a43ee57ff74ac8af01e/cvrp\_nls/utils.py\#L57-L86}{GFACS矩阵输入}。}
% content block 9:58 reference
\referencetext{[ER33] \href{https://github.com/ai4co/gfacs/blob/39d88cf2a5f9eb81eeea5a43ee57ff74ac8af01e/cvrp\_nls/test.py}{GFACS测试入口}。}
% content block 9:59 reference
\referencetext{[ER34] \href{https://proceedings.mlr.press/v306/fawkes26a.html}{f-Trajectory Balance}。}
% content block 9:60 reference
\referencetext{[ER35] \href{https://proceedings.mlr.press/v306/feng26o.html}{CAM}。}
% content block 9:61 reference
\referencetext{[ER36] \href{https://github.com/Optima-CityU/LLM4AD\_Next}{OpenLoopX LLM4AD Next工程}。}
% content block 9:62 reference
\referencetext{[ER37] \href{https://arxiv.org/html/2601.08308v1}{AgriAgent}。}
% content block 9:63 reference
\referencetext{[ER38] \href{https://arxiv.org/html/2606.18519v2\#S4.SS3}{As You Wish v2几何局限}。}
% content block 9:64 reference
\referencetext{[ER39] \href{https://ucmercedrobotics.github.io/asUwish.html}{As You Wish项目及schema}。}
% content block 9:65 reference
\referencetext{[ER40] \href{https://arxiv.org/pdf/2411.04983}{DINO-WM}。}
% content block 9:66 reference
\referencetext{[ER41] \href{https://arxiv.org/html/2512.24497v1}{JEPA-WMs}。}
% content block 9:67 reference
\referencetext{[ER42] \href{https://arxiv.org/html/2608.04653v1}{CoCo}。}
% content block 9:68 reference
\referencetext{[ER43] \href{https://github.com/unmannedlab/RELLIS-3D}{RELLIS-3D官方数据}。}
% content block 9:69 reference
\referencetext{[ER44] \href{https://theairlab.org/TartanDrive2/}{TartanDrive 2.0}。}
% content block 9:70 reference
\referencetext{[ER45] \href{https://github.com/castacks/tartan\_drive\_2.0}{TartanDrive 2.0下载工具}。}
% content block 9:71 reference
\referencetext{[ER46] \href{https://api.projectchrono.org/vehicle\_terrain.html}{Chrono地形模型}。}
% content block 9:72 reference
\referencetext{[ER47] \href{https://rl4co.ai4co.org/docs/content/api/envs/base/}{RL4CO环境接口}。}
% content block 9:73 reference
\referencetext{[ER48] \href{https://github.com/automl/DACBench}{DACBench}。}
% content block 9:74 reference
\referencetext{[ER49] \href{https://github.com/automl/AutoFolio}{AutoFolio}。}
% content block 9:75 reference
\referencetext{[ER50] \href{https://mlopez-ibanez.github.io/irace/}{irace}。}
% content block 9:76 reference
\referencetext{[ER51] \href{https://github.com/ahottung/NDS}{NDS前次调研代码入口}。}
% content block 9:77 reference
\referencetext{[ER52] \href{https://arxiv.org/abs/2501.03715}{NDS前次调研论文入口}。}
\end{document}
